\chapter{Probabilistic Foundations}
\label{cml:chap:probabilistic_foundations}

\begin{tldr}
This chapter supplies the vocabulary the rest of the volume leans on. Four ideas, each with a derivation an interviewer can ask for on a whiteboard. \textbf{Maximum likelihood} says every loss function is a noise model in disguise---Gaussian noise gives squared error, Laplace noise gives absolute error, Bernoulli gives cross-entropy---so choosing a loss \textit{is} choosing a generative story. \textbf{MAP} adds a prior and lands exactly on the regularizers of Chapter~\ref{cml:chap:supervised_core}: a Gaussian prior on the weights \textit{is} ridge with $\lambda = \sigma^2/\tau^2$, a Laplace prior \textit{is} lasso, and a Beta prior on a rate \textit{is} Laplace smoothing measured in pseudo-counts; \textbf{full Bayes} keeps the whole posterior, which matters for small samples, Thompson sampling (Section~\ref{cml:sec:bandits}), and any decision nonlinear in the parameter---and matters almost not at all once $n$ is large. The \textbf{bias--variance decomposition} is derived in [DL 2]; here it is \textit{instantiated}, with the exact algebra showing why ridge always beats OLS for some $\lambda>0$, why kNN's variance is exactly $\sigma^2/k$, and a knob table tying every hyperparameter in this volume to the term it moves. \textbf{Calibration} is the part production teams get wrong: a 0.92-AUC model can have useless probabilities, because ranking metrics are invariant to monotone score transforms and decision thresholds are not; the fix is understanding \textbf{proper scoring rules} and the calibration--refinement decomposition that explains why they reward honest probabilities. Finally, \textbf{uncertainty}: bootstrap mechanics and where they break, the pinball loss whose minimizer is provably the $\tau$-quantile, and split conformal prediction at 2026-interview depth---finite-sample marginal coverage with no distributional assumptions, and a precise account of everything that guarantee does \textit{not} give you.
\end{tldr}

%==============================================================================
\section{Maximum Likelihood}
\label{cml:sec:mle}
%==============================================================================

\subsection{The Likelihood, and Why It Is Not a Probability of the Parameters}

Fix a parametric family $p(y \mid x; \theta)$. Given data $\mathcal{D} = \{(x_i, y_i)\}_{i=1}^n$ assumed independent given the inputs, the \textbf{likelihood} is the probability the model assigns to the observed labels, read as a function of $\theta$ with the data held fixed:
\begin{equation}
L(\theta) = \prod_{i=1}^n p(y_i \mid x_i; \theta), \qquad \ell(\theta) = \log L(\theta) = \sum_{i=1}^n \log p(y_i \mid x_i; \theta).
\end{equation}
The maximum likelihood estimate is $\hat\theta_{\text{MLE}} = \argmax_\theta \ell(\theta)$. Two framing points that interviewers probe:

\begin{itemize}
    \item \textbf{The likelihood is not a distribution over $\theta$.} It does not integrate to one in $\theta$ and carries no probability semantics about the parameter---that is precisely what a prior buys you in Section~\ref{cml:sec:map}. Saying ``the likelihood tells us the probability that $\theta$ is correct'' is a red flag.
    \item \textbf{The log is not a convenience, it is the whole reason this is tractable.} Products of $10^6$ small numbers underflow; sums do not. And the log turns the exponential-family densities that dominate applied ML into the sums of simple terms whose gradients are clean (Chapter~\ref{cml:chap:supervised_core}).
\end{itemize}

\subsection{Two Worked MLEs Every Candidate Owes on Demand}

\begin{mathresult}{Bernoulli MLE: the Count Ratio}
With $y_i \in \{0,1\}$ i.i.d.\ Bernoulli($\theta$) and $k = \sum_i y_i$ successes in $n$ trials,
\begin{equation}
\ell(\theta) = k \log\theta + (n-k)\log(1-\theta), \qquad
\ell'(\theta) = \frac{k}{\theta} - \frac{n-k}{1-\theta} = 0 \;\Longrightarrow\; \hat\theta = \frac{k}{n}.
\end{equation}
Second derivative $\ell'' = -k/\theta^2 - (n-k)/(1-\theta)^2 < 0$, so this is the maximum. The estimator is unbiased, and it is also catastrophic at the boundary: $k=0$ gives $\hat\theta = 0$, i.e., ``this event is impossible,'' which is the failure a prior fixes (Section~\ref{cml:sec:map}).
\end{mathresult}

\begin{mathresult}{Gaussian MLE: the Sample Mean and the Biased Variance}
With $y_i \sim \mathcal{N}(\mu, \sigma^2)$ i.i.d.,
\begin{equation}
\ell(\mu,\sigma^2) = -\frac{n}{2}\log(2\pi\sigma^2) - \frac{1}{2\sigma^2}\sum_{i=1}^n (y_i - \mu)^2 .
\end{equation}
Setting $\partial\ell/\partial\mu = \sigma^{-2}\sum_i (y_i - \mu) = 0$ gives $\hat\mu = \bar{y}$. Substituting and setting $\partial\ell/\partial\sigma^2 = -\tfrac{n}{2\sigma^2} + \tfrac{1}{2\sigma^4}\sum_i (y_i-\bar y)^2 = 0$ gives
\begin{equation}
\hat\sigma^2_{\text{MLE}} = \frac{1}{n}\sum_{i=1}^n (y_i - \bar y)^2 .
\end{equation}
Because $\E\big[\sum_i (y_i - \bar y)^2\big] = (n-1)\sigma^2$, we get $\E[\hat\sigma^2_{\text{MLE}}] = \frac{n-1}{n}\sigma^2$: the MLE \textbf{underestimates the variance by a factor $(n-1)/n$}, a bias of $-\sigma^2/n$.
\end{mathresult}

\textbf{Why the variance MLE is biased, in one sentence:} it measures spread around $\bar y$, the point that \textit{minimizes} that spread by construction, rather than around the unknown true $\mu$; estimating the mean from the same data consumes one degree of freedom. Multiplying by $n/(n-1)$ (Bessel's correction) restores unbiasedness. The honest follow-up is that it rarely matters---the bias is $1/n$ and vanishes fast (at $n=1000$ it is a tenth of a percent), it only shows up when $n$ is genuinely small or when many small-$n$ variances are aggregated (per-user, per-merchant, per-experiment-cell), and the unbiased variance estimator is not the unbiased \textit{standard deviation} estimator anyway, since $\sqrt{\cdot}$ is concave and Jensen's inequality bites.

\subsection{Every Loss Is a Noise Model}

\begin{keyinsight}[Choosing a Loss Is Choosing a Likelihood]
Minimizing a loss $\sum_i \mathcal{L}(y_i, f(x_i))$ is maximum likelihood under the noise model $p(y \mid f) \propto e^{-\mathcal{L}(y, f)}$. This is the single most reusable idea in the chapter: when a problem hands you an unfamiliar target---counts, durations, proportions, censored survival times---you do not pick a loss from a menu, you write down the distribution the target plausibly follows and take its negative log-likelihood. Table~\ref{cml:tab:mle_losses} is the standard dictionary, and being able to derive any row of it on the spot is the ask.
\end{keyinsight}

\begin{table}[H]
\centering
\small
\caption{Noise model $\to$ negative log-likelihood $\to$ the loss you already use}
\label{cml:tab:mle_losses}
\begin{tabularx}{\textwidth}{lllX}
\toprule
\textbf{Target} & \textbf{Noise model} & \textbf{Resulting loss} & \textbf{Where it bites} \\
\midrule
Real & Gaussian, fixed $\sigma^2$ & Squared error (MSE) & The default; $\sigma^2$ drops out as a constant scale \\
Real & Gaussian, $\sigma^2(x)$ learned & $(y-\mu)^2/2\sigma^2 + \tfrac{1}{2}\log\sigma^2$ & Heteroscedastic regression; a variance head is aleatoric uncertainty \\
Real, outliers & Laplace (scale $b$) & Absolute error (MAE) & Heavier tails $\Rightarrow$ median, not mean; robust to outliers \\
Real, outliers & Student-$t$ & Robust/Huber-like & Huber is the practical stand-in with a quadratic core \\
Binary & Bernoulli & Binary cross-entropy & Logistic regression, GBDT log loss \\
$K$-class & Categorical & Cross-entropy & Softmax; gradient $p - y$ at the logits \\
Counts & Poisson ($\log\mu$ link) & Poisson deviance & Demand, claims, impressions; variance $=$ mean assumption \\
Positive, skewed & Gamma / Tweedie & Gamma deviance & Insurance severity, spend; multiplicative errors \\
\bottomrule
\end{tabularx}
\end{table}

\textbf{Derive one row to show the pattern.} Under $y_i = f(x_i) + \varepsilon_i$ with $\varepsilon_i \sim \mathcal{N}(0,\sigma^2)$,
\begin{equation}
-\log p(y_i \mid x_i) = \frac{(y_i - f(x_i))^2}{2\sigma^2} + \tfrac{1}{2}\log(2\pi\sigma^2),
\end{equation}
and the second term does not depend on $f$, so maximizing likelihood over $f$ \textit{is} minimizing squared error. Swap in the Laplace density $p(\varepsilon) = \tfrac{1}{2b}e^{-|\varepsilon|/b}$ and the same two lines give $|y_i - f(x_i)|/b$ plus a constant: absolute error. That is the whole trick, and it explains a fact candidates often recite without mechanism---MSE chases the conditional mean while MAE chases the conditional median---because those are exactly the point predictions the two likelihoods reward.

\begin{mathresult}{Cross-Entropy $=$ MLE $=$ Minimizing KL to the Empirical Distribution}
For a discrete target with empirical distribution $\hat p$ over the observed labels and model $q_\theta$,
\begin{equation}
\frac{1}{n}\sum_{i=1}^n -\log q_\theta(y_i) \;=\; H(\hat p, q_\theta) \;=\; \underbrace{H(\hat p)}_{\text{independent of }\theta} + \KL(\hat p \,\|\, q_\theta).
\end{equation}
So minimizing average cross-entropy is minimizing $\KL(\hat p \| q_\theta)$ over $\theta$, and the training loss can never reach zero: it is floored by $H(\hat p)$, the entropy of the labels themselves. Every improvement in log loss is literally bits saved against that floor---the framing that makes perplexity and the ``irreducible loss'' term in scaling laws [DL 2] mean something concrete.
\end{mathresult}

\subsection{What MLE Guarantees, and When It Falls Apart}

Under regularity conditions (identifiability, a smooth interior optimum, correct specification) the MLE is:
\begin{itemize}
    \item \textbf{Consistent}: $\hat\theta_n \to \theta_0$ in probability as $n \to \infty$.
    \item \textbf{Asymptotically normal and efficient}: $\sqrt{n}(\hat\theta_n - \theta_0) \to \mathcal{N}\big(0, I(\theta_0)^{-1}\big)$, where $I$ is the Fisher information per observation; the asymptotic variance attains the Cramér--Rao bound, so no other regular estimator beats it asymptotically. This is where standard errors and Wald confidence intervals in a logistic-regression summary table come from.
    \item \textbf{Invariant to reparameterization}: $\widehat{g(\theta)} = g(\hat\theta)$ for any function $g$. The MLE of the odds is the odds of the MLE---a property MAP does \textit{not} have (Section~\ref{cml:sec:map}).
\end{itemize}

\begin{warningbox}
All three properties are \textbf{asymptotic and conditional on correct specification}, and the interview value of naming them is knowing where they stop applying. Finite-sample MLE is biased (the Gaussian variance above). Under separation, the logistic likelihood has no finite maximizer and weights diverge (Chapter~\ref{cml:chap:supervised_core}). With $d$ comparable to $n$ the asymptotics are worthless. And when the model is misspecified---the usual case---the MLE converges to the parameter minimizing $\KL(p_{\text{true}} \| p_\theta)$, which is the closest wrong model, not the truth: consistency is consistency \textit{for the projection}, and sandwich (robust) standard errors exist precisely because the Fisher-information formula is wrong under misspecification.
\end{warningbox}

%==============================================================================
\section{MAP: Where Regularization Comes From}
\label{cml:sec:map}
%==============================================================================

\subsection{Posterior $\propto$ Likelihood $\times$ Prior}

Bayes' rule turns the likelihood into a distribution over parameters:
\begin{equation}
p(\theta \mid \mathcal{D}) = \frac{p(\mathcal{D} \mid \theta)\, p(\theta)}{p(\mathcal{D})} \;\propto\; \underbrace{p(\mathcal{D}\mid\theta)}_{\text{likelihood}} \cdot \underbrace{p(\theta)}_{\text{prior}} .
\end{equation}
The \textbf{maximum a posteriori} estimate is the mode of this posterior, and taking negative logs kills the normalizer:
\begin{equation}
\hat\theta_{\text{MAP}} = \argmax_\theta \big[\ell(\theta) + \log p(\theta)\big] = \argmin_\theta \big[\underbrace{-\ell(\theta)}_{\text{data-fit loss}} \; \underbrace{-\log p(\theta)}_{\text{penalty}}\big].
\end{equation}

\begin{keyinsight}[MAP Is Exactly Regularized MLE]
The penalty term in every regularized objective you have ever written is the negative log of a prior. \textbf{Regularization is not a trick; it is a prior belief about parameters, stated in log space.} The correspondence is exact and runs both ways: given a penalty you can read off the prior it encodes, and given a prior you can read off the penalty. The two canonical instances---Gaussian prior giving ridge and Laplace prior giving lasso---close the loop with the regularization geometry of Chapter~\ref{cml:chap:supervised_core}, and a candidate who can only offer ``L2 keeps weights small'' when probed for the probabilistic reading is signaling the shallow version of that chapter.
\end{keyinsight}

\subsection{Gaussian Prior $\Rightarrow$ Ridge}

\begin{mathresult}{Ridge Regression Is MAP with a Gaussian Prior, and $\lambda = \sigma^2/\tau^2$}
Take the likelihood $y_i \mid x_i \sim \mathcal{N}(x_i^\top\beta,\ \sigma^2)$ and the independent prior $\beta_j \sim \mathcal{N}(0, \tau^2)$. The negative log posterior is
\begin{align}
-\log p(\beta \mid \mathcal{D}) &= \frac{1}{2\sigma^2}\sum_{i=1}^n (y_i - x_i^\top\beta)^2 \;+\; \frac{1}{2\tau^2}\sum_{j=1}^d \beta_j^2 \;+\; \text{const} .
\end{align}
Multiplying by the positive constant $2\sigma^2$ does not move the argmin, giving
\begin{equation}
\hat\beta_{\text{MAP}} = \argmin_\beta \Big\{ \norm{y - X\beta}_2^2 + \lambda \norm{\beta}_2^2 \Big\}, \qquad \boxed{\lambda = \frac{\sigma^2}{\tau^2}} .
\end{equation}
\end{mathresult}

The formula $\lambda = \sigma^2/\tau^2$ is where the derivation stops being an algebra exercise and starts paying rent, because it says what $\lambda$ \textit{means}: the ratio of noise to prior signal. Noisy targets (large $\sigma^2$) or a tight belief that coefficients are small (small $\tau^2$) both argue for heavier shrinkage. Three consequences follow immediately, and stating them is what separates a recited derivation from an owned one:

\begin{itemize}
    \item \textbf{Why you standardize features first.} An isotropic prior $\mathcal{N}(0,\tau^2 I)$ says every coefficient is a priori the same size. That is a statement about units---halve a feature's scale and its coefficient doubles---so the prior is only coherent after standardization. This is the probabilistic reason behind the practical rule, not a separate fact.
    \item \textbf{Why you do not penalize the intercept.} There is no prior belief that the target's mean is near zero, so the intercept gets a flat prior and no penalty.
    \item \textbf{Why $\lambda$ should scale with noise.} Two datasets with the same features and different label noise want different $\lambda$; cross-validation discovers this empirically, and the formula explains what it discovered.
\end{itemize}

\subsection{Laplace Prior $\Rightarrow$ Lasso}

Replace the Gaussian prior with a Laplace (double-exponential) prior $p(\beta_j) = \frac{1}{2b}e^{-|\beta_j|/b}$. Then $-\log p(\beta) = \frac{1}{b}\sum_j |\beta_j| + \text{const}$, and the same rescaling gives
\begin{equation}
\hat\beta_{\text{MAP}} = \argmin_\beta \Big\{ \norm{y - X\beta}_2^2 + \lambda \norm{\beta}_1 \Big\}, \qquad \lambda = \frac{2\sigma^2}{b} .
\end{equation}
The Laplace density has a \textbf{sharp peak at zero and heavier tails} than a Gaussian of matched scale, which is the probabilistic statement of the geometric fact from Chapter~\ref{cml:chap:supervised_core}: it puts real prior mass on coefficients being \textit{exactly} zero-ish while remaining tolerant of a few genuinely large ones, whereas the Gaussian prior's smooth peak makes exact zeros a measure-zero event and its light tails punish large coefficients hard.

\begin{warningbox}
The lasso is the posterior \textbf{mode} under a Laplace prior---and the posterior \textbf{mean} under the same prior is not sparse at all. Every coefficient's posterior has continuous density, so its mean is almost surely nonzero. The sparsity that makes lasso useful is an artifact of reporting the mode of a spiky posterior, not a property of the Bayesian model. This is the cleanest concrete demonstration that MAP is not Bayesian inference, and it is a standard L7 probe. If you actually want Bayesian sparsity, you need a prior with mass \textit{at} zero (spike-and-slab) or a heavy-tailed shrinkage prior (horseshoe), not a Laplace.
\end{warningbox}

\begin{table}[H]
\centering
\small
\caption{The prior--penalty dictionary}
\label{cml:tab:prior_penalty}
\begin{tabularx}{\textwidth}{llX}
\toprule
\textbf{Prior on $\beta$} & \textbf{Penalty $-\log p(\beta)$} & \textbf{Reading} \\
\midrule
$\mathcal{N}(0,\tau^2)$ & $\lambda\norm{\beta}_2^2$, $\lambda = \sigma^2/\tau^2$ & Ridge; shrink everything proportionally, no exact zeros \\
Laplace$(0,b)$ & $\lambda\norm{\beta}_1$, $\lambda = 2\sigma^2/b$ & Lasso; sparse mode, dense posterior mean \\
Gaussian $\times$ Laplace & Elastic net & Grouped selection among correlated features \\
Flat / improper & None & MAP $=$ MLE; ``no prior'' is itself a prior choice \\
$\mathcal{N}(\beta_{\text{prev}}, \tau^2)$ & $\lambda\norm{\beta - \beta_{\text{prev}}}_2^2$ & Warm-start retraining: ``stay near yesterday's model'' \\
Beta$(\alpha,\beta)$ on a rate & Pseudo-counts $\alpha-1$, $\beta-1$ & Laplace smoothing; see below \\
Horseshoe / spike-and-slab & Non-convex & Sparsity without shrinking the large coefficients \\
\bottomrule
\end{tabularx}
\end{table}

The warm-start row is worth more than a glance: penalizing $\norm{\beta - \beta_{\text{prev}}}_2^2$ is a Gaussian prior centered on the previously deployed model, and it is exactly how production teams keep a nightly retrain from thrashing predictions on unchanged inputs. Written this way the prior variance $\tau^2$ becomes an explicit, arguable statement of how much the world is allowed to have changed since yesterday, rather than a hyperparameter someone tuned once.

\subsection{Priors as Pseudo-Counts: Beta-Bernoulli and Laplace Smoothing}

For rates---click-through, conversion, fraud---the prior's strength has an exact interpretation as \textit{data you pretend to have seen}. With a Beta$(\alpha,\beta)$ prior and $k$ successes in $n$ trials, the posterior is Beta$(\alpha + k,\ \beta + n - k)$ (derived in Section~\ref{cml:sec:bayes}), and its mode is
\begin{equation}
\hat\theta_{\text{MAP}} = \frac{\alpha + k - 1}{\alpha + \beta + n - 2},
\end{equation}
valid when $\alpha + k > 1$ and $\beta + n - k > 1$. Set $\alpha = \beta = 2$ and this becomes $\frac{k+1}{n+2}$: \textbf{add-one (Laplace) smoothing is exactly MAP under a Beta$(2,2)$ prior.} Add-$\kappa$ smoothing $\frac{k+\kappa}{n+2\kappa}$ is MAP under Beta$(\kappa+1, \kappa+1)$. The same $\frac{k+1}{n+2}$ is also the posterior \textit{mean} under the uniform Beta$(1,1)$ prior---a coincidence worth knowing, because it is the source of endless confusion about which prior ``is'' add-one smoothing; the honest answer is that it depends on whether you report the mode or the mean.

This is why the naive-Bayes smoothing constant of Chapter~\ref{cml:chap:supervised_core} is not a hack: it is a prior with a stated strength of $\alpha+\beta$ pseudo-observations, and it repairs the boundary catastrophe of the Bernoulli MLE---a word never seen in the training set gets $\hat\theta = 0$, which annihilates the entire product of a naive-Bayes posterior no matter what the other evidence says.

\subsection{Two Things MAP Is Not}

\textbf{MAP is not Bayesian inference.} It is a point estimate that throws away the posterior, and therefore every uncertainty-aware decision the posterior would have supported. The failure is not academic: in high dimensions the mode is systematically unrepresentative of where the posterior mass actually lives. For a $d$-dimensional standard Gaussian the mode sits at the origin, but essentially all the probability mass lies in a thin shell at radius $\approx \sqrt{d}$---the mode is in a region the distribution almost never visits.

\textbf{MAP is not invariant to reparameterization.} Densities transform with a Jacobian factor, so the mode moves under a nonlinear change of variables: the MAP estimate of $\theta$ and the MAP estimate of $\log\theta$ are not related by $\log$. A uniform prior on a probability is not uniform on the log-odds, so ``I used an uninformative prior'' is only meaningful relative to a chosen parameterization. The MLE \textit{is} invariant (Section~\ref{cml:sec:mle}), and so are the posterior mean under the correct transformation rules and the full posterior itself. Naming this asymmetry unprompted is a reliable senior signal.

%==============================================================================
\section{Full Bayes: Posteriors, Conjugacy, and When They Matter}
\label{cml:sec:bayes}
%==============================================================================

\subsection{The Posterior Predictive}

Full Bayesian inference keeps $p(\theta \mid \mathcal{D})$ and integrates it out when predicting, rather than plugging in a single $\hat\theta$:
\begin{equation}
p(\tilde y \mid \tilde x, \mathcal{D}) = \int p(\tilde y \mid \tilde x, \theta)\, p(\theta \mid \mathcal{D})\, d\theta
\quad\text{versus the plug-in}\quad p(\tilde y \mid \tilde x, \hat\theta).
\end{equation}
The difference between these two is the entire practical content of ``being Bayesian.'' The posterior predictive is \textbf{wider}, because it adds parameter uncertainty (epistemic) on top of observation noise (aleatoric); the plug-in prediction reports only the latter and is therefore overconfident, by an amount that shrinks as $O(1/n)$.

\subsection{Conjugacy, and Two Worked Examples}

A prior is \textbf{conjugate} to a likelihood when the posterior stays in the prior's family, so the integral above has a closed form and inference is arithmetic on the parameters. Conjugacy is a convenience, not a principle---it constrains what you can believe a priori---but the two conjugate pairs below cover a striking share of what industry Bayesian reasoning actually is.

\begin{mathresult}{Beta-Bernoulli}
Prior $\theta \sim \text{Beta}(\alpha,\beta)$, likelihood $k$ successes in $n$ Bernoulli trials. Since $p(\theta) \propto \theta^{\alpha-1}(1-\theta)^{\beta-1}$ and $p(\mathcal{D}\mid\theta) \propto \theta^{k}(1-\theta)^{n-k}$,
\begin{equation}
p(\theta\mid\mathcal{D}) \propto \theta^{\alpha+k-1}(1-\theta)^{\beta+n-k-1} \;\Longrightarrow\; \theta \mid \mathcal{D} \sim \text{Beta}(\alpha+k,\ \beta+n-k).
\end{equation}
The posterior mean is a \textbf{convex combination of the prior mean and the MLE}, with weights given by the prior's pseudo-sample-size $\alpha+\beta$ against the real sample size $n$:
\begin{equation}
\E[\theta \mid \mathcal{D}] = \frac{\alpha+k}{\alpha+\beta+n} = \underbrace{\frac{\alpha+\beta}{\alpha+\beta+n}}_{\text{weight on prior}}\cdot\frac{\alpha}{\alpha+\beta} \;+\; \underbrace{\frac{n}{\alpha+\beta+n}}_{\text{weight on data}}\cdot\frac{k}{n}.
\end{equation}
\end{mathresult}

\textbf{Worked numbers.} A new ad creative has 2 clicks in 20 impressions. The MLE says its CTR is $2/20 = 10\%$, which would make it the best creative in the account. Site-wide CTR is 3\%, and you encode that with Beta$(3, 97)$---prior mean $3/100 = 3\%$, prior strength 100 pseudo-impressions. The posterior is Beta$(5, 115)$, with
\begin{equation*}
\text{mean} = \frac{5}{120} = 4.17\%, \qquad \text{mode} = \frac{5-1}{120-2} = 3.39\%, \qquad \text{sd} = \sqrt{\tfrac{5 \cdot 115}{120^2 \cdot 121}} = 0.0182 ,
\end{equation*}
i.e.\ a posterior standard deviation of 1.82 percentage points. Check the shrinkage identity: $\tfrac{100}{120}(0.03) + \tfrac{20}{120}(0.10) = 0.025 + 0.0167 = 4.17\%$. \checkmark\ Three things fall out of this one calculation. The estimate moved from an implausible 10\% to a defensible 4.2\%. The posterior standard deviation of 1.8 points says the creative is not distinguishable from site average yet---the honest report is ``no evidence,'' not ``best creative.'' And the prior strength $\alpha+\beta = 100$ is a dial you can defend to a stakeholder in units they understand: ``we are treating this as if it had already run 100 impressions at the site average.'' Estimating $(\alpha,\beta)$ from the observed spread of CTRs across all creatives, rather than picking them, is \textbf{empirical Bayes}, and it is how this becomes a production system rather than a thought experiment.

\begin{mathresult}{Gaussian-Gaussian (known variance)}
Prior $\mu \sim \mathcal{N}(\mu_0, \tau_0^2)$, data $y_1,\dots,y_n \sim \mathcal{N}(\mu, \sigma^2)$ with $\sigma^2$ known. The posterior is $\mathcal{N}(\mu_n, \tau_n^2)$ where \textbf{precisions add} and the mean is a precision-weighted average:
\begin{equation}
\frac{1}{\tau_n^2} = \frac{1}{\tau_0^2} + \frac{n}{\sigma^2}, \qquad
\mu_n = \tau_n^2\left(\frac{\mu_0}{\tau_0^2} + \frac{n\bar y}{\sigma^2}\right).
\end{equation}
The posterior predictive for a new observation is
\begin{equation}
\tilde y \mid \mathcal{D} \sim \mathcal{N}\big(\mu_n,\ \underbrace{\tau_n^2}_{\text{epistemic}} + \underbrace{\sigma^2}_{\text{aleatoric}}\big),
\end{equation}
which is the cleanest formal statement of the confidence-interval-versus-prediction-interval distinction in Section~\ref{cml:sec:uncertainty}: $\tau_n^2 \to 0$ as $n$ grows, $\sigma^2$ never does.
\end{mathresult}

\subsection{When the Posterior Actually Matters}

\begin{keyinsight}[The Point Estimate Is Usually Enough---Until It Is Not]
\textbf{Bernstein--von Mises} says that under regularity conditions the posterior concentrates as $\mathcal{N}(\hat\theta_{\text{MLE}},\, I(\theta_0)^{-1}/n)$: with enough data the prior washes out, and MLE, MAP, and posterior mean agree to within $O(1/n)$. For a model fit on ten million rows, the parameter posterior is effectively a point mass and arguing about it is theater---the uncertainty that matters there is aleatoric noise and model misspecification, neither of which the parameter posterior describes. The posterior earns its compute in exactly four situations.
\end{keyinsight}

\begin{enumerate}
    \item \textbf{Small $n$, or small $n$ \textit{per unit}.} Aggregate data volume is irrelevant if the decision is per-merchant, per-creative, or per-arm and each unit has twenty observations. This is the single most common real case, and hierarchical/empirical-Bayes shrinkage across units is the standard answer.
    \item \textbf{The decision is nonlinear in the parameter.} Expected utility is $\E_\theta[u(\theta)]$, not $u(\E[\theta])$; whenever $u$ is convex or concave, Jensen's inequality makes the plug-in answer systematically wrong. Thresholds, caps, auctions, and inventory decisions are all nonlinear.
    \item \textbf{The algorithm consumes a posterior directly.} Thompson sampling draws $\theta_a$ from each arm's posterior and plays the argmax (Section~\ref{cml:sec:bandits}); with a Beta-Bernoulli posterior this is four lines of code and its exploration is automatically proportional to the probability that an arm is best. There is no point-estimate version of this algorithm.
    \item \textbf{You need probability statements about hypotheses.} ``$P(\text{B beats A}) = 0.93$'' is a posterior quantity and cannot be read off a p-value (Section~\ref{cml:sec:testing_rebuilt}); Bayesian A/B analysis and expected-loss stopping rules are the honest way to get it, with the same caveat about priors that frequentist analysis has about stopping rules.
\end{enumerate}

\textbf{The honest 2026 framing to say out loud:} full Bayes is rarely the production path for a supervised model---MCMC does not fit a nightly retrain and variational approximations trade the guarantee for tractability. But the \textit{reasoning} is load-bearing in bandits, experimentation, small-sample estimation, and anywhere a rate must be reported for a unit with almost no data. Claiming you would ``go fully Bayesian'' on a 100M-row CTR model is a red flag; failing to reach for a posterior when asked to rank 50,000 merchants by fraud rate with a median of 12 transactions each is a bigger one.

%==============================================================================
\section{Bias and Variance, Instantiated}
\label{cml:sec:bias_variance_classical}
%==============================================================================

The decomposition itself---the add-and-subtract expansion, the vanishing cross-term, and the caveat that the classical U-curve is a statement about the underparameterized regime only---is derived and owned in \textbf{[DL 2]}, together with double descent and scaling laws. Do not re-derive it here and do not re-derive it in an interview when the interviewer has already been given it; what this chapter adds is the part that makes the identity useful on classical models: \textit{which knob moves which term, and by exactly how much}.

\begin{mathresult}{The Identity, and What Each Expectation Is Over}
For squared loss with $y = f(x) + \varepsilon$, $\E[\varepsilon]=0$, $\Var(\varepsilon)=\sigma^2$, and $\hat f$ trained on a random dataset $\mathcal{D}$, at a fixed test point $x$:
\begin{equation}
\E_{\mathcal{D},\varepsilon}\big[(y - \hat f(x))^2\big] = \underbrace{\big(\E_\mathcal{D}[\hat f(x)] - f(x)\big)^2}_{\text{Bias}^2} + \underbrace{\Var_\mathcal{D}\big(\hat f(x)\big)}_{\text{Variance}} + \underbrace{\sigma^2}_{\text{irreducible}}
\end{equation}
\textbf{The expectations are over two independent sources of randomness}: the draw of the \textit{training set} (which is what $\E_\mathcal{D}$ and $\Var_\mathcal{D}$ average over---bias and variance are properties of a \textit{learning procedure}, not of a fitted model) and the \textit{noise in the test label} (which contributes $\sigma^2$). Being unable to name these when asked is the single most reliable red flag on this question, and it is the one interviewers plant the follow-up for.
\end{mathresult}

\subsection{The Irreducible Term, Stated Honestly}

$\sigma^2$ is the variance of $y$ given $x$: the part of the target no function of your features can explain. Three honest qualifications, all of which get asked:

\begin{itemize}
    \item \textbf{It is irreducible given your representation, not in nature.} $\sigma^2 = \Var(y \mid x)$ is conditional on the feature set you chose. Add a feature that carries some of that variation---device fingerprint, a merchant history aggregate, the text of the review---and $\sigma^2$ falls. ``Irreducible'' means ``no model can beat it,'' not ``no data scientist can beat it,'' and better features are the only knob in Table~\ref{cml:tab:bv_knobs} that moves this term at all.
    \item \textbf{You cannot separate it from bias with one dataset.} Empirically you observe test error, which is $\text{Bias}^2 + \text{Var} + \sigma^2$; you can estimate variance by refitting on resamples, but the remaining lump is bias plus noise and no amount of held-out data splits it. You need replicate labels at the same $x$ (repeated measurements, multiple annotators) or a trusted lower bound---human agreement rates, the Bayes error of a known-optimal reference---to pin $\sigma^2$ down. This is exactly why annotator-agreement studies are worth their cost.
    \item \textbf{Label noise is often not irreducible, it is a data-quality bug.} A large apparent $\sigma^2$ frequently turns out to be inconsistent labeling guidelines, delayed-outcome mislabeling (a chargeback that arrives after the label snapshot), or a leaky join. Diagnose before you accept.
\end{itemize}

\subsection{Ridge: Buying Variance with Bias, Priced Exactly}

Ridge is the cleanest place to see the trade as algebra rather than a slogan. Take an orthonormal design ($X^\top X = I_d$, so $\hat\beta^{\text{OLS}}$ has covariance $\sigma^2 I$). Then $\hat\beta^{\text{ridge}} = \hat\beta^{\text{OLS}}/(1+\lambda)$, and for a single coefficient:
\begin{equation}
\text{Bias}(\hat\beta_j^{\text{ridge}}) = \frac{\beta_j}{1+\lambda} - \beta_j = -\frac{\lambda\beta_j}{1+\lambda}, \qquad
\Var(\hat\beta_j^{\text{ridge}}) = \frac{\sigma^2}{(1+\lambda)^2},
\end{equation}
\begin{equation}
\text{MSE}(\lambda) = \frac{\lambda^2\beta_j^2 + \sigma^2}{(1+\lambda)^2}, \qquad
\frac{d\,\text{MSE}}{d\lambda} = \frac{2(\lambda\beta_j^2 - \sigma^2)}{(1+\lambda)^3} .
\end{equation}

\begin{keyinsight}[There Is Always a $\lambda > 0$ That Beats OLS]
At $\lambda = 0$ the derivative is $-2\sigma^2 < 0$: MSE is \textit{strictly decreasing} as you leave the unbiased estimator. Whenever the labels carry any noise at all, some positive amount of shrinkage improves squared-error risk---unbiasedness is not a virtue worth paying for. The optimum is at
\begin{equation}
\lambda^\star = \frac{\sigma^2}{\beta_j^2},
\end{equation}
which is \textbf{the same formula as the MAP result $\lambda = \sigma^2/\tau^2$} from Section~\ref{cml:sec:map}, with the prior variance $\tau^2$ playing the role of the typical squared coefficient. The Bayesian derivation and the frequentist risk calculation agree on the answer, which is the strongest possible way to end this question.
\end{keyinsight}

\textbf{Numbers, because interviewers ask for them.} With $\sigma^2 = 1$ and $\beta_j = 2$, we get $\lambda^\star = 1/4$. Then $\text{bias}^2 = \frac{(0.25)^2 \cdot 4}{(1.25)^2} = 0.16$ and $\Var = \frac{1}{(1.25)^2} = 0.64$, for MSE $= 0.80$ against OLS's $0 + 1.00 = 1.00$: variance fell 36\%, bias rose from zero to 0.16, and the net is a 20\% risk reduction. Push to $\lambda = 1$ and it reverses---bias$^2 = 1.00$, $\Var = 0.25$, MSE $= 1.25$, worse than OLS. That is the U-curve, computed rather than gestured at.

\subsection{kNN: The Cleanest Knob in Classical ML}

For kNN regression at $x_0$, averaging the $k$ nearest training targets,
\begin{equation}
\E\big[(y_0 - \hat f(x_0))^2\big] = \underbrace{\sigma^2}_{\text{irreducible}} + \underbrace{\Big[f(x_0) - \tfrac{1}{k}\textstyle\sum_{l=1}^{k} f(x_{(l)})\Big]^2}_{\text{Bias}^2} + \underbrace{\frac{\sigma^2}{k}}_{\text{Variance}} .
\end{equation}
The variance term is \textbf{exactly $\sigma^2/k$}---averaging $k$ independent noise draws---so it falls monotonically in $k$ with no approximation. The bias term is the gap between $f(x_0)$ and the average of $f$ over the neighborhood, which grows as $k$ pulls in more distant, less relevant neighbors. At $k=1$: zero training error, variance $\sigma^2$, and a prediction that changes entirely if one label flips. At $k=n$: the global mean, variance $\sigma^2/n$, and a bias equal to the full deviation of $f(x_0)$ from the average. This is the decomposition with both terms in closed form and the knob visible in each, which is why it is the standard teaching example and a fair thing to be asked to reproduce.

\subsection{Ensembles: Two Ensembles, Two Terms}

Chapter~\ref{cml:chap:trees_ensembles} derives the mechanics; the decomposition is the language that organizes them.

\textbf{Bagging and random forests attack variance and leave bias alone.} The trees are identically distributed, so $\E[\frac{1}{B}\sum_b T_b] = \E[T_1]$---averaging cannot move the expectation, hence cannot move bias. It moves the variance to $\rho\sigma^2 + \frac{1-\rho}{B}\sigma^2$ (Section~\ref{cml:sec:bagging_rf}), which is why RF wants \textit{deep} trees (bias is the term the ensemble cannot fix, so the base learner must not have any) and why adding trees cannot overfit (variance is monotone decreasing in $B$, bias is constant). Feature subsampling attacks $\rho$, the correlation floor that caps how much variance averaging can remove.

\textbf{Boosting attacks bias and inflates variance as it goes.} Each round fits the negative gradient of the loss at the current predictions, systematically correcting what the ensemble is still getting wrong---bias falls with $M$. But the ensemble is being fit sequentially to the same data, so its variance grows with $M$ and eventually it fits label noise: \textbf{boosting overfits in rounds, which is why early stopping is not optional} (Section~\ref{cml:sec:boosting_theory}). This is why boosting wants \textit{shallow} trees---a high-bias base learner is exactly what a bias-reducing procedure needs, and depth is an interaction-order budget, not a capacity dial.

That asymmetry---which ensemble can be run too long, and which cannot---is the sharpest instantiation of the decomposition in classical ML and the reason ``more trees overfits a random forest'' is a reliable red flag.

\begin{table}[H]
\centering
\small
\caption{The knob table: which hyperparameter moves which term. ``$\uparrow$'' means increasing the knob.}
\label{cml:tab:bv_knobs}
\begin{tabularx}{\textwidth}{llllX}
\toprule
\textbf{Model} & \textbf{Knob} & \textbf{Bias} & \textbf{Var} & \textbf{Mechanism} \\
\midrule
kNN & $k \uparrow$ & $\uparrow$ & $\downarrow$ & Variance is exactly $\sigma^2/k$; neighbors get farther \\
Ridge/lasso & $\lambda \uparrow$ & $\uparrow$ & $\downarrow$ & Shrinkage toward the prior; optimum at $\sigma^2/\tau^2$ \\
Linear model & Add features & $\downarrow$ & $\uparrow$ & Richer hypothesis class, more parameters to estimate \\
Poly/spline & Degree $\uparrow$ & $\downarrow$ & $\uparrow$ & The textbook U-curve \\
SVM & $C \uparrow$, $\gamma \uparrow$ & $\downarrow$ & $\uparrow$ & Harder margin, narrower kernel, wigglier boundary \\
Tree & \texttt{max\_depth} $\uparrow$ & $\downarrow$ & $\uparrow$ & Finer partition, fewer points per leaf \\
Tree & \texttt{min\_samples\_leaf} $\uparrow$ & $\uparrow$ & $\downarrow$ & Coarser leaves, more averaging per prediction \\
RF & $B \uparrow$ & --- & $\downarrow$ & I.i.d.\ trees: expectation fixed, variance $\to \rho\sigma^2$ \\
RF & \texttt{mtry} $\downarrow$ & $\uparrow$ & $\downarrow$ & Lowers $\rho$ (the floor) at some per-tree cost \\
GBM & Rounds $M \uparrow$ & $\downarrow$ & $\uparrow$ & Sequential bias correction; overfits without early stopping \\
GBM & $\eta \downarrow$ (with $M \uparrow$) & --- & $\downarrow$ & Smaller steps at matched fit: the shrinkage effect \\
GBM & Tree depth $\uparrow$ & $\downarrow$ & $\uparrow$ & Higher interaction order per round \\
Any & $n \uparrow$ & --- & $\downarrow$ & More data shrinks variance; bias and $\sigma^2$ untouched \\
Any & Better features & $\downarrow$ & --- & \textbf{The only row that moves $\sigma^2$}, by moving signal out of the noise \\
\bottomrule
\end{tabularx}
\end{table}

\subsection{Diagnosis, and Where the Decomposition Stops}

\textbf{Diagnosis in practice} is the train-versus-validation gap read alongside the learning curve. High training \textit{and} validation error, converging as $n$ grows to a plateau above your target: high bias---add capacity, add features, reduce regularization, raise depth or rounds. Low training error with a large validation gap that narrows as $n$ grows: high variance---more data, more regularization, more averaging, fewer features. The learning curve is the discriminator, because it says whether more data will help: if the two curves have already met, more rows will not move anything and the problem is bias or the feature set.

\textbf{Where it stops.} The clean three-term identity is a squared-loss fact. Under 0--1 loss there is no additive decomposition into non-negative bias and variance terms: the published unified decompositions (Domingos, 2000; James, 2003) carry an interaction term that can be \textbf{negative}, meaning variance can \textit{reduce} classification error---when the average prediction sits on the wrong side of the boundary, dispersion around it lets some draws land on the right side. This is a real effect and part of why a heavily biased model like naive Bayes can classify well while producing nonsense probabilities (Section~\ref{cml:sec:calibration}). And for the modern regime---the interpolation threshold, double descent, why the U-curve does not describe overparameterized networks---go to \textbf{[DL 2]}; the classical decomposition remains the right mental model for tabular models with a real train--test gap, which is essentially everything in this volume.

%==============================================================================
\section{Calibration and Proper Scoring Rules}
\label{cml:sec:calibration}
%==============================================================================

\subsection{What Calibration Is, and Three Things It Is Not}

A scoring model is \textbf{calibrated} if, among all cases where it predicts $p$, the event happens a fraction $p$ of the time:
\begin{equation}
\prob{Y = 1 \mid \hat p(X) = p} = p \qquad \text{for all } p \text{ in the model's range}.
\end{equation}
Three clarifications carry most of the interview value:

\begin{itemize}
    \item \textbf{Calibration is not accuracy, and it is not hard to achieve on its own.} The constant predictor $\hat p(x) \equiv \pi$, where $\pi$ is the base rate, is \textit{perfectly calibrated} and completely useless. Calibration is a necessary condition for a probability to be meaningful, never a sufficient one for the model to be good. What you want is calibration \textit{plus} \textbf{sharpness} (predictions spread out toward 0 and 1) or, equivalently, calibration plus \textbf{refinement}---and Section~\ref{cml:sec:proper_scores} shows that a proper scoring rule measures exactly that combination.
    \item \textbf{Calibration is not per-prediction correctness.} It is an aggregate property of a population of predictions. A single prediction of 0.7 is never right or wrong; only the set of all 0.7-predictions can be.
    \item \textbf{Marginal calibration does not imply subgroup calibration.} A model can be perfectly calibrated overall while being systematically overconfident on one segment and underconfident on another, with the errors cancelling in the aggregate. This is why calibration must be checked on the slices you actually make decisions about---per region, per merchant tier, per device---and it is the technical core of several fairness criteria.
\end{itemize}

\textbf{Measurement.} The \textbf{reliability diagram}---bin predictions, plot mean predicted probability against observed frequency, compare to the diagonal---and the \textbf{expected calibration error} that summarizes it into a scalar are the standard instruments; the formula and the temperature/Platt/isotonic fixes are cataloged in \textbf{[DL 23]}. What this chapter owns is why that scalar is more fragile than it looks, and what to reach for instead.

\subsection{ECE Is a Binning Artifact With a Number Attached}

\begin{warningbox}
\textbf{ECE is not a property of the model; it is a property of the model, the binning scheme, and the sample size.} Three specific failures to be able to name:

\textbf{1. It is upward-biased at finite sample.} Even a \textit{perfectly calibrated} model shows nonzero ECE, because the empirical frequency in a bin of $m$ points fluctuates around the true probability. For a bin whose true probability is $p$, the expected absolute deviation is approximately $\sqrt{2p(1-p)/(\pi m)}$. At $p = 0.5$ and $m = 100$ points per bin that is $0.040$; at $m = 1000$ it falls to $0.013$. \textbf{So a perfectly calibrated model evaluated with 15 bins of 100 points each reports ECE $\approx 0.04$}, and ``we improved ECE from 0.05 to 0.04'' may be measuring nothing. More bins means fewer points per bin means more upward bias---the estimator gets worse exactly when you try to make it finer.

\textbf{2. Equal-width and equal-mass binning give different answers.} On a fraud model where 99\% of scores are below 0.05, equal-width bins put nearly all the data in the first bin and the reported ECE is dominated by whatever happens there, while the sparse high-score bins---the ones that drive every decline decision---contribute almost nothing. Equal-mass (quantile) bins fix the sample-size imbalance but move the bin edges as the score distribution drifts, so ECE is no longer comparable across model versions.

\textbf{3. ECE is not a proper scoring rule and can be gamed.} A model can lower its binned ECE by making its outputs \textit{less} informative---snapping predictions to bin centers, or shrinking everything toward the base rate---without becoming a better probability estimator. Optimizing ECE directly is therefore a bad idea; use it as a diagnostic alongside a reliability diagram, and report a proper score as the headline number.
\end{warningbox}

The practical protocol: report a proper score (log loss or Brier) as the metric, a reliability diagram as the diagnostic, and ECE---with the binning scheme stated and held fixed---only as a communication device. Where a single number really is needed, prefer a debiased or kernel-smoothed calibration estimate over raw binned ECE, and always compute it on the decision-relevant slices.

\subsection{Proper Scoring Rules, and Why They Decompose}
\label{cml:sec:proper_scores}

\begin{mathresult}{Definition}
A scoring rule $S(p, y)$ (a loss, lower is better) is \textbf{proper} if, when the true conditional probability is $q$, the expected score
$\E_{y \sim q}[S(p,y)]$ is minimized at $p = q$; it is \textbf{strictly proper} if $p = q$ is the \textit{unique} minimizer. A strictly proper rule is a scoring system under which honest reporting is the uniquely optimal strategy---the model has no way to score better by lying about its beliefs.
\end{mathresult}

\textbf{Log loss is strictly proper.} With true probability $q$ and report $p$,
\begin{equation}
L(p) = -q\log p - (1-q)\log(1-p), \qquad L'(p) = -\frac{q}{p} + \frac{1-q}{1-p} = \frac{p-q}{p(1-p)} .
\end{equation}
This is zero only at $p = q$, negative for $p<q$ and positive for $p>q$, and $L''(p) = q/p^2 + (1-q)/(1-p)^2 > 0$, so $p=q$ is the unique minimum. The one-line version that is worth memorizing because it also gives the decomposition:
\begin{equation}
\boxed{\;L(p) - L(q) = q\log\frac{q}{p} + (1-q)\log\frac{1-q}{1-p} = \KL(q \,\|\, p) \;\ge\; 0\;}
\end{equation}
with equality if and only if $p = q$. \textbf{The excess log loss you pay for reporting the wrong probability is exactly the KL divergence from the truth to your report.}

\textbf{Brier is strictly proper.} With $B(p) = \E_{y\sim q}[(p-y)^2] = q(1-p)^2 + (1-q)p^2$,
\begin{equation}
B'(p) = 2(p - q), \qquad B(p) - B(q) = (p-q)^2 \ge 0 ,
\end{equation}
so again $p=q$ uniquely. The two rules differ in how they punish confident mistakes: Brier's penalty is bounded (at most 1), log loss's is unbounded ($-\log p \to \infty$ as $p \to 0$). That makes log loss the right objective when confidently-wrong predictions are catastrophic and the wrong one when a handful of mislabeled rows can dominate your entire evaluation---a real problem on datasets with label noise, where a single $p = 10^{-6}$ on a mislabeled positive contributes 13.8 nats.

\begin{mathresult}{The Calibration--Refinement Decomposition}
Group predictions by their value. Let $\bar y_p = \prob{Y=1 \mid \hat p = p}$ be the true frequency among cases where the model said $p$. Then for log loss,
\begin{equation}
\E[\text{log loss}] = \underbrace{\E_p\big[\KL(\bar y_p \,\|\, p)\big]}_{\text{calibration (}\ge 0\text{, zero iff calibrated)}} \;+\; \underbrace{\E_p\big[H(\bar y_p)\big]}_{\text{refinement (sharpness)}} ,
\end{equation}
which follows by applying the boxed identity within each group. The Brier score admits Murphy's (1973) three-term form,
\begin{equation}
\text{BS} = \underbrace{\tfrac{1}{n}\textstyle\sum_k n_k(p_k - \bar y_k)^2}_{\text{reliability (calibration)}} - \underbrace{\tfrac{1}{n}\textstyle\sum_k n_k(\bar y_k - \bar y)^2}_{\text{resolution}} + \underbrace{\bar y(1-\bar y)}_{\text{uncertainty (data, not model)}} ,
\end{equation}
with refinement $=$ uncertainty $-$ resolution, recovering the same two-term split. \textbf{This is why proper scoring rules incentivize calibration}: the calibration term is non-negative, is zero exactly when the model is calibrated, and is \textit{additively separate} from the refinement term---so you can never buy a better score by being miscalibrated, only by being sharper.
\end{mathresult}

\textbf{A worked check, four predictions.} Predict $(0.5, 0.5, 0.9, 0.9)$ with labels $(1,0,1,1)$. Directly, $\text{BS} = \frac{0.25+0.25+0.01+0.01}{4} = 0.13$. Now decompose: the $p=0.5$ group has $\bar y = 0.5$, the $p=0.9$ group has $\bar y = 1.0$, and the overall base rate is $\bar y = 0.75$. Reliability $= \frac{2(0)^2 + 2(0.1)^2}{4} = 0.005$; resolution $= \frac{2(0.25)^2 + 2(0.25)^2}{4} = 0.0625$; uncertainty $= 0.75 \times 0.25 = 0.1875$. Then $0.005 - 0.0625 + 0.1875 = 0.13$. \checkmark\ The same data under log loss: calibration $= \tfrac{1}{2}\KL(1.0\|0.9) = 0.0527$ nats, refinement $= \tfrac{1}{2}H(0.5) = 0.3466$ nats, total $0.3993$, which matches the direct computation $-\tfrac{1}{4}[2\log 0.5 + 2\log 0.9] = 0.3993$. \checkmark\ Note what the decomposition reveals: this model's Brier score is dominated by the \textit{uncertainty} of the data ($0.1875$), not by its own miscalibration ($0.005$)---which is why raw Brier scores are not comparable across datasets with different base rates.

\textbf{Accuracy and AUC are not strictly proper.} Accuracy at a 0.5 threshold has expected loss $1-q$ if you report anything above 0.5 and $q$ otherwise; reporting $p = q$ is \textit{a} minimizer, but so is every other value on the same side of the threshold. It is proper but not \textit{strictly} proper, which is exactly the failure that matters: it gives you no incentive to distinguish $0.51$ from $0.99$. AUC is worse in a more interesting way---it is invariant to any strictly increasing transformation of the scores, so it literally cannot see the difference between $p$ and $\sigma(10\,\mathrm{logit}\,p)$. That invariance is the property that makes it a good discrimination metric and a useless calibration metric, and it is the whole content of ``0.92 AUC, useless probabilities.''

\subsection{Which Classical Models Are Calibrated, and Why Not}

\begin{table}[H]
\centering
\small
\caption{The classical calibration map. ``Out of the box'' means before any post-hoc calibration.}
\label{cml:tab:calib_map}
\begin{tabularx}{\textwidth}{llX}
\toprule
\textbf{Model} & \textbf{Out of the box} & \textbf{Mechanism} \\
\midrule
Logistic regression & Good & Trained on a strictly proper score with the matching link; the score equations force the residuals to sum to zero, so mean prediction $=$ base rate \textit{in sample} \\
Naive Bayes & Badly overconfident & Conditional-independence violation double-counts correlated evidence; posteriors slam to 0/1 while the \textit{ranking} stays useful \\
SVM & No probabilities at all & Hinge loss produces a margin, not a likelihood; distances are not probabilities. Platt scaling was invented for exactly this \\
Random forest & Under-confident, compressed & The vote fraction averages $\{0,1\}$ votes and is almost never unanimous, so predictions bunch away from 0 and 1 \\
GBDT (log loss) & Fair, drifting overconfident & Proper training loss helps, but many rounds push margins outward; check after early stopping \\
AdaBoost & Badly miscalibrated & Exponential loss is proper for a \textit{different} link; scores need a sigmoid map before they mean anything \\
kNN & Coarse & Probabilities are $j/k$: only $k+1$ distinct values, and $k=10$ can never express 0.05 \\
Neural networks & Overconfident & See [DL 23] (Guo et al.); temperature scaling is the default fix there \\
Downsampled negatives & Systematically shifted & Any model trained on a rebalanced base rate; correct with the odds formula below \\
\bottomrule
\end{tabularx}
\end{table}

Two rows deserve their own paragraph. \textbf{Logistic regression's calibration is a theorem with fine print}: including an intercept, the stationarity condition $X^\top(y - \hat p) = 0$ makes the average predicted probability equal the observed base rate and forces zero residual correlation with every feature---so it is calibrated in-sample along every direction the features span. That guarantee does \textit{not} survive heavy regularization (which biases $\hat p$ toward the prior), a misspecified link, or distribution shift between training and serving. ``Logistic regression is always calibrated'' is a red flag; ``logistic regression is calibrated in-sample by construction, which is why it is the fair baseline for calibration comparisons'' is the right version.

\textbf{Negative downsampling is the most common production miscalibration and the easiest to fix.} If you keep negatives with probability $w$ and train on the resulting sample, the model's output $p_s$ relates to the true probability $p$ by an odds rescaling, and inverting gives
\begin{equation}
p = \frac{w\,p_s}{w\,p_s + (1 - p_s)} .
\end{equation}
Worked: a true 1\% CTR with $w = 0.1$ presents as $0.01/(0.01 + 0.99 \times 0.1) = 9.17\%$ in the resampled data; feeding $p_s = 0.0917$ and $w = 0.1$ back through the formula recovers $0.0100$. \checkmark\ Ranking is unaffected (the map is monotone), which is exactly why teams ship this bug for months---the AUC dashboard stays green while every downstream expected-value computation is off by an order of magnitude.

\subsection{Platt Scaling vs.\ Isotonic Regression}

Both fit a monotone map from raw scores to probabilities on a \textbf{held-out} calibration set. The choice is a bias--variance decision on the calibration map itself, which is a satisfying place for this chapter to land.

\textbf{Platt scaling} fits $\hat p = \sigma(a\,s + b)$---two parameters, estimated by logistic regression on (score, label) pairs. It works with a few hundred calibration points and is stable, but it can only express a sigmoid-shaped distortion; if the true miscalibration has a different shape (say a random forest's compression at both ends, which is roughly the \textit{inverse} of a sigmoid), Platt leaves residual error no amount of data will remove. Platt's original recipe also replaces the hard 0/1 targets with $\frac{N_+ + 1}{N_+ + 2}$ and $\frac{1}{N_- + 2}$, a Beta$(2,2)$ prior on the calibration labels---the same pseudo-count idea as Section~\ref{cml:sec:map}, applied one level up.

\textbf{Isotonic regression} fits the best non-decreasing step function, using the pool-adjacent-violators algorithm. It is nonparametric and can correct any monotone distortion, but its effective parameter count grows with the number of blocks, so it \textbf{overfits hard on small calibration sets}---below roughly 1{,}000 points it will chase noise and produce a step function that generalizes worse than the two-parameter fit. It also emits piecewise-constant outputs: every score inside a block collapses to a single value, creating ties that can destroy fine-grained ranking in exactly the top-of-list region an auction or a top-$k$ retrieval cares about. Being strictly increasing, Platt preserves AUC exactly; isotonic is only non-decreasing, so it can shave AUC slightly by creating those ties.

\textbf{The operating rule:} under $\sim$1{,}000 calibration points use Platt (or temperature scaling for a neural net, [DL 23]); above $\sim$10{,}000, isotonic is usually safe and strictly more expressive; in between, choose by cross-validated log loss on the calibration split. And in all cases the calibration set must be \textbf{disjoint from training}: calibrating on training-set predictions is the classic mistake, and it is worst for exactly the models that most need calibrating, because a tree ensemble's in-sample predictions are already near-perfect and the fitted map is then close to the identity.

\begin{keyinsight}[Ranking Metrics Do Not Care; Decisions Do]
AUC, NDCG, MRR, and every other rank-based metric are invariant to strictly monotone transforms of the score, so they are structurally incapable of detecting miscalibration. Any use of the score as a \textit{number} is not: an expected-value computation ($\text{bid} = \text{pCTR} \times \text{value}$), a budget pacing controller, a cascade routing rule, or a fixed threshold. Decision theory gives the threshold explicitly---with cost $c_{\text{FP}}$ for a false positive and $c_{\text{FN}}$ for a false negative, acting is optimal when $(1-p)c_{\text{FP}} < p\,c_{\text{FN}}$, i.e.
\begin{equation*}
p^\star = \frac{c_{\text{FP}}}{c_{\text{FP}} + c_{\text{FN}}} .
\end{equation*}
For fraud with a \$5 cost to decline a good transaction and a \$120 chargeback, $p^\star = 5/125 = 0.04$---not 0.5, and not 0.9. \textbf{A threshold is only meaningful if the scores are calibrated and the number was derived from costs}; a team running at ``0.9'' usually chose it by staring at a precision--recall table, which silently re-derives a cost ratio nobody wrote down.
\end{keyinsight}

%==============================================================================
\section{Uncertainty Quantification}
\label{cml:sec:uncertainty}
%==============================================================================

\subsection{Aleatoric vs.\ Epistemic}

\textbf{Aleatoric} uncertainty is randomness in the outcome given everything you know: two transactions with identical features, one fraudulent and one not; the fair coin; the delivery that is late because of traffic your features do not observe. \textbf{Epistemic} uncertainty is uncertainty about the model itself, arising from finite data: had you trained on a different sample, you would have predicted something different.

The operative distinction is what shrinks. \textbf{Epistemic uncertainty shrinks with more data} and vanishes in the limit; \textbf{aleatoric does not}---and ``more data fixes aleatoric noise'' is a canonical red flag. The honest qualification, the same one as the irreducible-error term in Section~\ref{cml:sec:bias_variance_classical}, is that aleatoric uncertainty is defined \textit{relative to your feature set}: better features convert aleatoric into explained variation. More rows never will; more columns might.

Mapping to the decomposition: $\sigma^2$ is aleatoric, $\Var_\mathcal{D}(\hat f)$ is epistemic, and bias is neither---it is misspecification, and it is the term that no uncertainty estimate in this section will warn you about. Practical proxies: bootstrap or ensemble spread estimates epistemic uncertainty; residual spread at a densely-sampled $x$, or a learned variance head (Table~\ref{cml:tab:mle_losses}), estimates aleatoric; a Bayesian posterior predictive gives both at once, since its variance is $\tau_n^2 + \sigma^2$.

\subsection{Confidence Interval vs.\ Prediction Interval}

\begin{keyinsight}[The Distinction That Gets Asked Constantly]
A \textbf{confidence interval} covers a \textit{parameter} or a \textit{conditional mean}: ``the average delivery time for this route is 31--33 minutes.'' A \textbf{prediction interval} covers a \textit{new observation}: ``the next delivery on this route takes 18--46 minutes.'' In OLS at a test point $x_0$,
\begin{equation*}
\text{CI: } \hat y_0 \pm t\,\hat\sigma\sqrt{x_0^\top (X^\top X)^{-1} x_0},
\qquad
\text{PI: } \hat y_0 \pm t\,\hat\sigma\sqrt{1 + x_0^\top (X^\top X)^{-1} x_0} .
\end{equation*}
The entire difference is the \textbf{$1$} under the second radical: the new observation's own noise. As $n \to \infty$ the CI width goes to zero while the PI width converges to $2t\hat\sigma \approx 3.92\hat\sigma$. \textbf{A prediction interval never shrinks below the noise floor}, no matter how much data you collect.
\end{keyinsight}

\textbf{Numbers.} With $\sigma = 10$ minutes and $n = 400$, the 95\% CI half-width for the mean is $1.96 \times 10/\sqrt{400} = 0.98$ minutes, while the PI half-width is $1.96 \times 10\sqrt{1 + 1/400} = 19.6$ minutes---twenty times wider. Quoting a CI when a stakeholder asked ``how late might \textit{my} package be'' understates the answer by that factor, and it is the mistake behind a great many overconfident product promises.

\subsection{The Bootstrap}

\textbf{Mechanics.} To get a sampling distribution for any statistic $\hat\theta$ with no formula: draw $B$ resamples of size $n$ \textit{with replacement} from the data, recompute $\hat\theta^{*b}$ on each, and treat $\{\hat\theta^{*b}\}$ as a stand-in for the sampling distribution. The percentile interval takes the $[\alpha/2, 1-\alpha/2]$ quantiles; the basic (reverse-percentile) interval $[2\hat\theta - \theta^*_{1-\alpha/2},\ 2\hat\theta - \theta^*_{\alpha/2}]$ corrects for the bootstrap distribution's bias, and BCa additionally corrects skew. Rule of thumb on $B$: a few hundred resamples suffice for a standard error, but $B \ge 2000$ for percentile intervals, because you are estimating tail quantiles and the tails are where resamples are scarce.

The logic is a substitution: the bootstrap approximates ``resample from the population'' by ``resample from the empirical distribution.'' Everything that can go wrong follows from that one substitution being a bad approximation.

\begin{warningbox}
\textbf{When the bootstrap fails---name at least three.}
\begin{itemize}
    \item \textbf{Extremes.} The bootstrap distribution of the sample maximum is degenerate: no resample can exceed the observed max, and the max is missing entirely from about 36.8\% of resamples (each point appears with probability $1 - (1-1/n)^n \approx 0.632$---the same constant as out-of-bag sampling in Section~\ref{cml:sec:bagging_rf}). The same problem afflicts extreme quantiles and any statistic driven by a handful of points.
    \item \textbf{Dependence.} I.i.d.\ resampling destroys autocorrelation and clustering, producing intervals that are far too narrow. For time series use a block bootstrap (Chapter~\ref{cml:chap:time_series}); for grouped data---users with many sessions, patients with many visits---resample the \textit{clusters}, not the rows. Bootstrapping rows when the dependence is at the user level is the same error as the wrong randomization unit in Section~\ref{cml:sec:metric_design}.
    \item \textbf{Tiny $n$.} With $n = 10$ the empirical distribution is a poor stand-in for the population and coverage is unreliable; the bootstrap is asymptotic machinery wearing nonparametric clothes.
    \item \textbf{Heavy tails and non-smooth functionals.} The bootstrap for the mean is inconsistent when the variance is infinite, and it is not consistent for statistics that are non-smooth functions of the data, including post-selection quantities like ``the coefficient of whichever variable lasso chose.''
    \item \textbf{Resampling the wrong step.} If feature selection, hyperparameter search, or threshold tuning happened before you started resampling, the interval is too narrow because it never sees that variability. The whole pipeline goes inside the bootstrap loop---the same principle as the cross-validation hygiene of Chapter~\ref{cml:chap:applied_craft}.
\end{itemize}
\end{warningbox}

\subsection{Quantile Regression and the Pinball Loss}

If you need the spread of the outcome rather than its mean, model the quantiles directly. The \textbf{pinball} (check) loss at level $\tau \in (0,1)$, with residual $u = y - q$, is
\begin{equation}
\rho_\tau(u) = u\big(\tau - \mathds{1}[u < 0]\big) =
\begin{cases}
\tau\, u & u \ge 0 \quad (\text{under-prediction}) \\
(1-\tau)|u| & u < 0 \quad (\text{over-prediction}).
\end{cases}
\end{equation}
It is an asymmetrically weighted absolute error: at $\tau = 0.9$, being too low costs nine times as much per unit as being too high, which pushes the fit upward until exactly 10\% of the mass sits above it.

\begin{mathresult}{The Pinball Minimizer Is the $\tau$-Quantile}
Let $Y$ have CDF $F$. The expected loss of a constant prediction $q$ is
\begin{equation}
R(q) = \tau\!\int_q^\infty (y-q)\,dF(y) \;+\; (1-\tau)\!\int_{-\infty}^q (q-y)\,dF(y) .
\end{equation}
Differentiating under the integral (the boundary terms vanish because the integrands are zero at $y=q$),
\begin{equation}
R'(q) = -\tau\big(1 - F(q)\big) + (1-\tau)F(q) = F(q) - \tau .
\end{equation}
Setting $R'(q) = 0$ gives $F(q) = \tau$, so $q^\star = F^{-1}(\tau)$. Since $R''(q) = f(q) \ge 0$, $R$ is convex and this is the global minimum. Sanity check: $\tau = 0.5$ gives $\rho = \tfrac{1}{2}|u|$, whose minimizer is the median---absolute error and the median are the special case, exactly as squared error and the mean are for the Gaussian.
\end{mathresult}

Fitting $\rho_\tau$ conditionally on $x$ (linear quantile regression, or a GBDT with the pinball objective---the quantile row of Table~\ref{cml:tab:boosting_losses}) gives conditional quantiles, and fitting $\tau = 0.05$ and $\tau = 0.95$ gives a nominal 90\% prediction interval whose width \textit{varies with $x$}, which is what you want: uncertainty should be larger in sparse regions. Two practical caveats: separately fitted quantile models can \textbf{cross} ($\hat q_{0.05}(x) > \hat q_{0.95}(x)$ for some $x$), fixed by joint estimation, post-hoc sorting, or monotonicity constraints; and the coverage is only as good as the model---\textbf{there is no finite-sample guarantee that the nominal 90\% interval covers 90\%}. That gap is precisely what conformal prediction closes.

\subsection{Conformal Prediction}
\label{cml:sec:conformal}

Conformal prediction converts \textit{any} point predictor into a set predictor with a finite-sample coverage guarantee, assuming only exchangeability. It has moved from a niche literature to a standard interview name-drop, and the follow-ups are now about the fine print rather than the recipe.

\begin{mathresult}{Split (Inductive) Conformal Prediction}
\textbf{Recipe.}
\begin{enumerate}
    \item Split the data into a proper training set and a \textbf{calibration set} of size $n$, disjoint from training.
    \item Fit $\hat f$ on the training set only. The model is now frozen; conformal never touches it again.
    \item Compute \textbf{nonconformity scores} on the calibration set. For regression, $s_i = |y_i - \hat f(x_i)|$; for classification, $s_i = 1 - \hat p_{y_i}(x_i)$, the softmax mass \textit{not} placed on the true label.
    \item Take $\hat q$ = the $\big\lceil (n+1)(1-\alpha)\big\rceil / n$ empirical quantile of $\{s_1,\dots,s_n\}$---note the $(n+1)$, which is the finite-sample correction that makes the guarantee exact rather than asymptotic.
    \item Predict the set $C(x_{n+1}) = \{y : s(x_{n+1}, y) \le \hat q\}$. For regression this is the fixed-width interval $\hat f(x_{n+1}) \pm \hat q$; for classification it is $\{y : \hat p_y(x_{n+1}) \ge 1 - \hat q\}$, a set whose size varies with the input.
\end{enumerate}
\textbf{Guarantee.} If the calibration points and the test point are exchangeable,
\begin{equation}
1 - \alpha \;\le\; \prob{Y_{n+1} \in C(X_{n+1})} \;\le\; 1 - \alpha + \frac{1}{n+1} .
\end{equation}
\textbf{Proof idea, in one line:} exchangeability makes the rank of $s_{n+1}$ among $s_1,\dots,s_{n+1}$ uniform on $\{1,\dots,n+1\}$, so $\prob{s_{n+1} \le s_{(k)}} \ge k/(n+1)$; take $k = \lceil (n+1)(1-\alpha)\rceil$. That is the entire argument, and being able to give it is what distinguishes knowing conformal prediction from having heard of it.
\end{mathresult}

\textbf{Worked mechanics.} With $n = 1000$ calibration points and $\alpha = 0.1$: $\lceil 1001 \times 0.9 \rceil = 901$, so take the 901st smallest absolute residual---the 90.1th percentile, not the 90th. Feasibility requires $\lceil (n+1)(1-\alpha)\rceil \le n$, i.e.\ $n \ge 1/\alpha - 1$: you cannot get 90\% coverage from 8 calibration points, and any implementation must handle that case (the correct output is the infinite interval, not a crash).

\begin{warningbox}
\textbf{What the guarantee does \textit{not} promise.} Four items, in the order interviewers ask them.

\textbf{1. It is marginal, not conditional.} $\prob{Y \in C(X)} \ge 1-\alpha$ averages over the distribution of $X$. It does not say $\prob{Y \in C(X) \mid X = x} \ge 1-\alpha$, and distribution-free conditional coverage is provably impossible with non-trivial sets. In practice, 90\% marginal coverage routinely means 98\% on easy inputs and 65\% on a hard subgroup. If you need per-group coverage, calibrate \textit{within} each group (Mondrian/group-conditional conformal); if you want widths that adapt to input difficulty, use conformalized quantile regression, which conformalizes a quantile-regression interval and inherits both its adaptivity and the coverage guarantee.

\textbf{2. It says nothing about width.} A model that outputs the base rate for every input still achieves exact coverage---with enormous intervals. \textbf{Coverage is validity; width is efficiency, and only the underlying model buys efficiency.} Conformal is a wrapper, not an improvement.

\textbf{3. The coverage is also marginal over the calibration draw.} For a single fixed calibration set, the realized coverage is random: it follows $\text{Beta}(n+1-\ell,\ \ell)$ with $\ell = \lfloor (n+1)\alpha \rfloor$. At $n = 100$, $\alpha = 0.1$, that is Beta$(91,10)$: mean 0.901, \textbf{standard deviation 3.0 points}, so your particular deployment might really be covering 85\% or 95\%. At $n = 1000$ the standard deviation drops to 0.95 points. Size the calibration set for the coverage precision you need, not just for the guarantee.

\textbf{4. Exchangeability is the whole ballgame, and distribution shift voids it.} Under temporal drift, feedback loops, or any train/serve mismatch, calibration residuals are no longer exchangeable with the test point and the guarantee is simply gone---silently, with no error message. The repairs, all worth naming: \textit{weighted conformal} restores validity under covariate shift \textit{if} you know or can estimate the likelihood ratio $w(x) = dP_{\text{test}}/dP_{\text{train}}$; \textit{adaptive conformal inference} adjusts $\alpha_t$ online in response to observed coverage, trading the finite-sample guarantee for a long-run-average one and is the right default for streaming/time-series (Chapter~\ref{cml:chap:time_series}); and the cheapest practical mitigation is to recalibrate on a rolling recent window and \textbf{monitor realized coverage as a production metric}, which turns a broken assumption into an alert instead of a surprise.
\end{warningbox}

\begin{table}[H]
\centering
\small
\caption{Choosing an uncertainty tool}
\label{cml:tab:uncertainty_tools}
\begin{tabularx}{\textwidth}{llX}
\toprule
\textbf{Tool} & \textbf{Captures} & \textbf{Guarantee, and what limits it} \\
\midrule
Bootstrap & Epistemic & Asymptotic only. Needs i.i.d.\ data; fails on extremes, on dependence, and at tiny $n$ \\
Bayesian posterior & Both & Exact \textit{given the model}. Needs a prior, a likelihood, and compute; wrong under misspecification \\
Quantile regression & Aleatoric & None. Coverage is only as good as the model; separately fitted quantiles can cross \\
Split conformal & Marginal coverage & Finite-sample and distribution-free. Marginal only; voided by shift; width comes from the base model \\
Ensemble spread & Epistemic (proxy) & None. Correlated members understate it; cheap, model-agnostic, widely used \\
Replicate labels & Aleatoric & Exact. Requires repeated measurements at the same $x$, which most datasets lack \\
\bottomrule
\end{tabularx}
\end{table}

The composite answer a strong candidate gives to ``how would you put uncertainty on this model'': separate the two kinds, use a bootstrap or ensemble for epistemic and a quantile head or variance head for aleatoric, wrap the whole thing in split conformal so the reported interval has a coverage guarantee rather than a hope, and monitor realized coverage in production because the exchangeability assumption is the first thing drift breaks.

%==============================================================================
\section{Interview Questions}
\label{cml:sec:probabilistic_interview}
%==============================================================================

\begin{interviewq}
\textbf{Q: MLE, MAP, and fully Bayesian---define each, and tell me where regularization comes from.}

\badgeLfive{L5} \quad \badgetype{First Principles} \quad \badgefreq{\freqhigh}

\stronganswer

They are three points on one axis: how much of the posterior you keep. \textbf{MLE} maximizes $\ell(\theta) = \sum_i \log p(y_i \mid x_i;\theta)$---no prior, and the likelihood is not a distribution over $\theta$. \textbf{MAP} maximizes $\ell(\theta) + \log p(\theta)$, the mode of the posterior. \textbf{Full Bayes} keeps $p(\theta\mid\mathcal{D})$ and integrates it out when predicting, giving a posterior predictive rather than a plug-in.

Regularization falls out of the middle one immediately. Negating the MAP objective gives $-\ell(\theta) - \log p(\theta)$: a data-fit loss plus a penalty, where \textbf{the penalty is exactly the negative log prior}. Instantiate it: a Gaussian likelihood with an i.i.d.\ $\mathcal{N}(0,\tau^2)$ prior on the weights gives $\norm{y-X\beta}^2 + \lambda\norm{\beta}_2^2$ with $\lambda = \sigma^2/\tau^2$---ridge. Swap in a Laplace prior and you get the $\ell_1$ penalty---lasso, with $\lambda = 2\sigma^2/b$. On a rate, a Beta$(2,2)$ prior gives $(k+1)/(n+2)$, which is add-one smoothing. So ``regularization strength'' is literally the ratio of noise variance to prior variance, and ``prior strength'' on a rate is literally a count of pseudo-observations.

When to use each: MLE when data is plentiful relative to parameters and you have no defensible prior; MAP whenever you would have regularized anyway---which is almost always---and knowing the prior interpretation tells you what $\lambda$ means and why you standardize features first and skip the intercept; full Bayes when $n$ is small \textit{per unit of decision} (per merchant, per arm, per creative), when the decision is nonlinear in the parameter so $\E[u(\theta)] \ne u(\E[\theta])$, when the algorithm literally consumes a posterior (Thompson sampling), or when you owe someone a probability statement about a hypothesis. Close with the honest bound: Bernstein--von Mises says the posterior concentrates as $\mathcal{N}(\hat\theta_{\text{MLE}}, I^{-1}/n)$, so at large $n$ all three agree and the distinction is theater---the uncertainty that matters at scale is aleatoric noise and misspecification, which the parameter posterior does not describe.

\redflags
\begin{itemize}
    \item ``MAP is Bayesian inference.'' It is a point estimate that discards the posterior---and in high dimensions the mode is not even where the mass lives.
    \item ``Regularization is just a trick to prevent overfitting,'' with no prior reading available when probed. This is the specific follow-up the question exists for.
    \item Calling the likelihood ``the probability that $\theta$ is right.''
    \item Proposing full Bayes for a 100M-row production model as if compute were free, or refusing to ever use a prior on principle.
\end{itemize}

\interviewtest{Whether the regularizers you use daily are organized in your head as probabilistic statements rather than as library arguments. The tell is whether $\lambda = \sigma^2/\tau^2$ arrives unprompted.}

\goodgreat
\begin{itemize}
    \item \badgeLfive{L5} Correct definitions, MAP $=$ regularized MLE, names the Gaussian$\to$L2 and Laplace$\to$L1 correspondences.
    \item \badgeLsix{L6} Derives $\lambda = \sigma^2/\tau^2$, uses it to explain feature standardization and the unpenalized intercept, and gives concrete regimes where the posterior earns its cost (small $n$ per unit, Thompson sampling, nonlinear decisions).
    \item \badgeLseven{L7} Notes that MAP is not reparameterization-invariant while MLE is, that the lasso's sparsity is an artifact of reporting a mode (the Laplace posterior mean is dense), and that Bernstein--von Mises makes the whole distinction vanish at scale.
\end{itemize}

\followups{``Why is a uniform prior not `no assumption'?'' ``What would a prior centered on last week's deployed model give you?'' ``Where does $\lambda$ come from if you cannot cross-validate?''}
\end{interviewq}

\begin{interviewq}
\textbf{Q: Show that ridge regression is MAP estimation. What is $\lambda$ in terms of the model?}

\badgeLfive{L5} \quad \badgetype{Mathematical} \quad \badgefreq{\freqmed}

\stronganswer

State the two ingredients, then turn the crank. Likelihood: $y_i \mid x_i \sim \mathcal{N}(x_i^\top\beta, \sigma^2)$. Prior: $\beta_j \sim \mathcal{N}(0,\tau^2)$ independently. Posterior $\propto$ likelihood $\times$ prior, so the negative log posterior is
\begin{equation*}
-\log p(\beta\mid\mathcal{D}) = \frac{1}{2\sigma^2}\sum_i (y_i - x_i^\top\beta)^2 + \frac{1}{2\tau^2}\norm{\beta}_2^2 + \text{const}.
\end{equation*}
Multiply by $2\sigma^2$---a positive constant, so the argmin is unchanged---and you have $\norm{y - X\beta}_2^2 + \lambda\norm{\beta}_2^2$ with $\lambda = \sigma^2/\tau^2$. Then say what the formula \textit{means}: $\lambda$ is the ratio of noise to prior signal. Noisier labels or a tighter prior both argue for more shrinkage, and this is why standardizing features first is not a separate convention---an isotropic prior asserts that all coefficients are a priori comparable in size, which is a claim about units and is only coherent after scaling---and why the intercept is unpenalized, since there is no prior belief that the target's mean is near zero.

The strong close is the frequentist half. In an orthonormal design, $\hat\beta^{\text{ridge}}_j = \hat\beta^{\text{OLS}}_j/(1+\lambda)$, so $\text{bias} = -\lambda\beta_j/(1+\lambda)$, $\Var = \sigma^2/(1+\lambda)^2$, and MSE $= (\lambda^2\beta_j^2 + \sigma^2)/(1+\lambda)^2$. Its derivative is $2(\lambda\beta_j^2 - \sigma^2)/(1+\lambda)^3$, which is negative at $\lambda=0$---so \textit{some} shrinkage always beats OLS---and vanishes at $\lambda^\star = \sigma^2/\beta_j^2$: the same formula the Bayesian derivation gave, with $\tau^2$ read as the typical squared coefficient. The two derivations agreeing is the best possible ending.

\redflags
\begin{itemize}
    \item Writing the posterior but forgetting that $\sigma^2$ appears in the likelihood, and therefore reporting $\lambda = 1/\tau^2$.
    \item Being unable to say what happens to $\lambda$ as $\tau \to \infty$ (it goes to zero: a flat prior recovers OLS) or as $\tau \to 0$ (all coefficients shrink to zero).
    \item Claiming ridge gives sparse solutions, or that unbiasedness is intrinsically desirable.
\end{itemize}

\interviewtest{Can you execute a three-line derivation and then interpret its output? Most candidates get the algebra; far fewer convert $\lambda = \sigma^2/\tau^2$ into the practical rules it implies.}

\goodgreat
\begin{itemize}
    \item \badgeLfive{L5} Clean derivation with the correct $\lambda = \sigma^2/\tau^2$.
    \item \badgeLsix{L6} Adds the standardization and intercept consequences, and does the Laplace-prior version for lasso with $\lambda = 2\sigma^2/b$.
    \item \badgeLseven{L7} Adds the risk calculation showing MSE strictly decreases at $\lambda = 0$ and is minimized at $\sigma^2/\beta_j^2$, and notes that the lasso is a posterior mode whose corresponding posterior mean is dense---so the Bayesian reading explains L1's sparsity as an artifact of mode-reporting, not a Bayesian property.
\end{itemize}

\followups{``Do the Laplace-prior version.'' ``What prior gives elastic net?'' ``What is the prior behind `stay close to last week's model' in a warm-started retrain?''}
\end{interviewq}

\begin{interviewq}
\textbf{Q: Why is the MLE of the Gaussian variance biased? Does it matter?}

\badgeLfive{L5} \quad \badgetype{Mathematical} \quad \badgefreq{\freqmed}

\stronganswer

Derive it first. With $\ell(\mu,\sigma^2) = -\tfrac{n}{2}\log(2\pi\sigma^2) - \tfrac{1}{2\sigma^2}\sum_i (y_i - \mu)^2$, the mean equation gives $\hat\mu = \bar y$ and the variance equation gives $\hat\sigma^2 = \tfrac{1}{n}\sum_i (y_i - \bar y)^2$. Since $\E\big[\sum_i (y_i - \bar y)^2\big] = (n-1)\sigma^2$, we get $\E[\hat\sigma^2] = \tfrac{n-1}{n}\sigma^2$---biased low by $\sigma^2/n$.

The mechanism in one sentence: the estimator measures spread around $\bar y$, which is by construction the point minimizing that spread in this sample, rather than around the unknown true $\mu$; estimating the mean from the same data consumes one degree of freedom. Multiplying by $n/(n-1)$---Bessel's correction---restores unbiasedness.

Does it matter? Usually not: the bias is $O(1/n)$, so at $n = 1000$ it is a tenth of a percent, swamped by everything else. It does matter in three places worth naming. When $n$ is genuinely small---per-cell variances in an experiment, per-user variances, small strata---where the correction is a several-percent effect on a quantity that then enters a denominator. When many small-$n$ variances are aggregated, since the bias does not average away. And in any variance-normalized statistic, where a systematically small denominator inflates the statistic. The sophisticated caveat: unbiasedness is not automatically what you want. The MLE has lower MSE than the unbiased estimator here, and $\sqrt{\hat\sigma^2}$ is a biased estimate of $\sigma$ regardless of which variance estimator you square-root, because the square root is concave and Jensen's inequality bites.

\redflags
\begin{itemize}
    \item ``Divide by $n-1$'' as a memorized rule with no account of why a degree of freedom is lost.
    \item Claiming the correction makes the \textit{standard deviation} unbiased.
    \item Treating unbiasedness as strictly better---the ridge question in this same chapter is the counterexample.
\end{itemize}

\interviewtest{Whether you can do a two-parameter MLE and reason about the estimator afterwards. It is also a proxy for statistical maturity: the ``does it matter'' half is the real question.}

\goodgreat
\begin{itemize}
    \item \badgeLfive{L5} Derives both MLEs, states the $(n-1)/n$ factor and the degrees-of-freedom explanation.
    \item \badgeLsix{L6} Quantifies when it matters (small $n$, aggregated small-$n$ estimates, variance-normalized statistics) and notes the MLE's lower MSE.
    \item \badgeLseven{L7} Connects to the general point that MLE is only asymptotically unbiased and asymptotically efficient, that all such guarantees are conditional on correct specification, and that under misspecification the MLE converges to the KL projection with sandwich standard errors rather than Fisher-information ones.
\end{itemize}

\followups{``What is the MLE of $\sigma$, and is it unbiased?'' ``What happens to these guarantees when the model is misspecified?'' ``How would a prior on $\sigma^2$ change the estimate?''}
\end{interviewq}

\begin{interviewq}
\textbf{Q: Training RMSE 0.31, validation RMSE 0.94, and the two learning curves are still far apart at 200K rows. Diagnose it and name the knobs, on the model of your choice.}

\badgeLfive{L5} \quad \badgetype{Debugging} \quad \badgefreq{\freqhigh}

\stronganswer

That gap with curves that have not met is the variance signature: the model is fitting sample-specific structure, and the still-converging curves say \textbf{more data will help}, which is the single most decision-relevant thing the plot tells you. Say what each term means first---bias is the error of the average model over training-set draws, variance is the spread across those draws, and $\sigma^2$ is the noise floor---and note that bias and variance are properties of the \textit{procedure}, not of the fitted object, because the expectation is over the training-set draw. The full identity and the deep-learning caveat live in [DL 2]; what matters here is which knob moves which term.

Then work the knobs by model. \textbf{GBDT}: reduce rounds via early stopping (boosting overfits in $M$ by construction), lower the learning rate and raise rounds, cut depth or \texttt{num\_leaves}, raise \texttt{min\_child\_weight}, add row/column subsampling, raise $\lambda$. \textbf{Random forest}: this pattern is less common because averaging is already a variance killer, so suspect leakage or a broken split first; then raise \texttt{min\_samples\_leaf} and lower \texttt{max\_features}. \textbf{Ridge/lasso}: raise $\lambda$, and note that $\lambda = \sigma^2/\tau^2$ means noisier targets want more of it. \textbf{kNN}: raise $k$---its variance is exactly $\sigma^2/k$, the cleanest knob in classical ML. Across all of them: more rows, fewer features, more aggressive averaging.

Then the seniority move---audit the measurement before spending a week on tuning. A gap this large invites two non-model explanations. \textbf{Leakage or a split defect}: a temporal split done randomly, group leakage where the same entity appears in both sides, or a target-derived feature will produce exactly this signature; check by comparing a random split against a temporal or grouped one. \textbf{Distribution shift between the two sets}: if validation is a later period, the gap may be drift rather than variance, and no regularization fixes drift. Only after those are cleared is ``tune for variance'' the right story.

\redflags
\begin{itemize}
    \item Reciting the decomposition but unable to say what the expectation is over---the specific probe interviewers plant here.
    \item Reading the numbers as high bias, or prescribing ``add capacity.''
    \item Jumping straight to hyperparameters without questioning the split.
    \item Claiming the trade-off means you can never reduce both---more data, better features, and ensembling all do.
\end{itemize}

\interviewtest{Whether the decomposition is a working diagnostic or a memorized formula. The knob mapping and the split audit are the two things that separate a candidate who has debugged models from one who has read about them.}

\goodgreat
\begin{itemize}
    \item \badgeLfive{L5} Correct diagnosis, correct direction on three or four knobs, names the learning-curve reading.
    \item \badgeLsix{L6} Names what each expectation averages over, gives the knob map across model families with mechanisms (variance $=\sigma^2/k$ for kNN, boosting overfits in $M$ while RF cannot overfit in $B$), and audits the split before tuning.
    \item \badgeLseven{L7} Separates $\sigma^2$ from bias honestly---you cannot split them from one dataset without replicate labels or a human-agreement bound---observes that only better features move the noise floor, and notes that the additive decomposition is a squared-loss fact that does not carry to 0--1 loss, where the interaction term can be negative.
\end{itemize}

\followups{``How would you estimate the irreducible noise?'' ``The curves have already met---now what?'' ``Does this analysis change for a 7B-parameter network?''}
\end{interviewq}

\begin{interviewq}
\textbf{Q: Our fraud model has 0.92 AUC, but the risk team says the probabilities are useless---their auto-decline rule fires at 0.9 and almost nothing crosses it. What is wrong, and how do you fix it?}

\badgeLsix{L6} \quad \badgetype{Debugging} \quad \badgefreq{\freqhigh}

\stronganswer

Separate discrimination from calibration in the first sentence. AUC measures \textit{ordering} and is invariant to any strictly increasing transform of the scores---it literally cannot distinguish $p$ from $\sigma(10\,\mathrm{logit}\,p)$---so 0.92 AUC says the model ranks fraud above non-fraud and says nothing whatsoever about whether a score of 0.9 means 90\% fraud. Their rule consumes the score as a \textit{number}; the metric they were shown does not evaluate it as one.

Diagnose in this order. \textbf{First, plot a reliability diagram} on a held-out set and report a proper score (log loss or Brier) alongside AUC; report ECE only with the binning scheme stated, since a perfectly calibrated model with 15 bins of 100 points still shows ECE $\approx 0.04$ from sampling noise alone. \textbf{Second, check the base rate.} With 0.3\% fraud, a well-calibrated model \textit{should} rarely exceed 0.9---the complaint may be that the threshold is wrong, not the model. \textbf{Third, check for negative downsampling}, the most common production cause: if negatives were kept with probability $w$, the scores live on a shifted base rate and must be mapped back with $p = w p_s / (w p_s + 1 - p_s)$. A true 1\% rate downsampled at $w = 0.1$ presents as 9.2\%---a factor of nine, invisible to AUC. \textbf{Fourth, check the model family}: a boosted ensemble pushes scores outward as rounds accumulate, a random forest compresses vote fractions away from 0 and 1, naive Bayes double-counts correlated evidence and slams to the extremes, and an SVM has no probabilities at all. \textbf{Fifth, check for drift}: calibration degrades under shift faster than AUC does, so a model calibrated at launch may be badly off now.

Then fix, in two independent parts. \textbf{Recalibrate}: fit a monotone map on a \textit{held-out} calibration set---Platt (two parameters) under about a thousand points, isotonic above ten thousand, chosen by cross-validated log loss in between. Both are monotone, so AUC is essentially unchanged; isotonic can shave it slightly by creating ties, which matters if anything downstream ranks within the plateau. \textbf{Re-derive the threshold from costs} rather than from a round number: acting is optimal when $p > c_{\text{FP}}/(c_{\text{FP}} + c_{\text{FN}})$, so a \$5 cost to decline a good transaction against a \$120 chargeback implies $p^\star = 0.04$, not 0.9. Finally, make it durable: monitor calibration (not only AUC) as a production metric, recalibrate on a rolling window, and never calibrate on training predictions---for a tree ensemble those are near-perfect and the fitted map comes out close to the identity.

\redflags
\begin{itemize}
    \item ``AUC is 0.92, the model is fine''---the exact conflation the question is built to catch.
    \item Proposing to retrain the model from scratch when a two-parameter post-hoc map on held-out data is the right first move.
    \item Fitting the calibrator on training-set predictions.
    \item Not asking about the base rate or about downsampling---the two cheapest explanations.
    \item Treating 0.9 as a given rather than asking where the number came from.
\end{itemize}

\interviewtest{Whether you know that ranking metrics and decision metrics measure different things, whether you can enumerate causes by prior probability instead of guessing one, and whether you treat the threshold as a decision-theoretic quantity rather than a constant someone typed in.}

\goodgreat
\begin{itemize}
    \item \badgeLfive{L5} Distinguishes discrimination from calibration, proposes a reliability diagram and a post-hoc calibrator.
    \item \badgeLsix{L6} Enumerates causes in likelihood order, catches downsampling with the correction formula, picks Platt vs.\ isotonic by calibration-set size, and re-derives the threshold from the cost ratio.
    \item \badgeLseven{L7} Reframes the metric contract---any consumer of the score as a number needs a proper scoring rule in the dashboard, not AUC---adds slice-wise calibration (marginal calibration does not imply subgroup calibration), monitors realized calibration under drift, and notes that ECE itself is binning-dependent and gameable so the headline number should be a proper score.
\end{itemize}

\followups{``Does isotonic regression change the AUC?'' ``How much calibration data do you need?'' ``The model is calibrated overall but wrong on one merchant segment---what now?'' ``Where did the 0.9 come from?''}
\end{interviewq}

\begin{interviewq}
\textbf{Q: Prove that log loss is a proper scoring rule. Why should I care?}

\badgeLsix{L6} \quad \badgetype{Mathematical} \quad \badgefreq{\freqmed}

\stronganswer

Define properness first: a scoring rule $S(p,y)$ is proper if, when the true probability is $q$, the expected score $\E_{y\sim q}[S(p,y)]$ is minimized at $p = q$, and strictly proper if that minimizer is unique. Then compute. For log loss,
\begin{equation*}
L(p) = \E_{y\sim q}[-\log p_y] = -q\log p - (1-q)\log(1-p).
\end{equation*}
Differentiate: $L'(p) = -q/p + (1-q)/(1-p) = (p-q)/\big(p(1-p)\big)$, which is negative for $p<q$, positive for $p>q$, and zero only at $p = q$; and $L''(p) = q/p^2 + (1-q)/(1-p)^2 > 0$, so the stationary point is the unique global minimum. Strictly proper.

The slicker version, which also gives you the payoff, is to subtract the minimum:
\begin{equation*}
L(p) - L(q) = q\log\frac{q}{p} + (1-q)\log\frac{1-q}{1-p} = \KL(q\,\|\,p) \ge 0,
\end{equation*}
with equality iff $p = q$ by Gibbs' inequality. \textbf{The excess log loss from reporting the wrong probability is exactly the KL divergence from the truth to your report.} The same two lines do Brier: $B(p) = q(1-p)^2 + (1-q)p^2$ gives $B'(p) = 2(p-q)$ and $B(p) - B(q) = (p-q)^2$.

Why care: because it is what makes a metric \textit{incentive-compatible}. Under a strictly proper rule the model's uniquely optimal strategy is to report its true beliefs---there is no way to score better by shading probabilities. Group predictions by their value and the score splits as
\begin{equation*}
\E[\text{log loss}] = \underbrace{\E_p[\KL(\bar y_p \| p)]}_{\text{calibration, }\ge 0} + \underbrace{\E_p[H(\bar y_p)]}_{\text{refinement}},
\end{equation*}
using the boxed identity within each group; Brier has Murphy's equivalent reliability $-$ resolution $+$ uncertainty form. The calibration term is non-negative, is zero exactly when the model is calibrated, and is \textit{additively separate} from refinement---so optimizing a proper score cannot be gamed by miscalibration and can only be improved by becoming sharper. Contrast: accuracy at a fixed threshold is proper but \textit{not strictly} proper---any report on the correct side of the threshold ties, so it gives no incentive to distinguish 0.51 from 0.99---and AUC is invariant to monotone transforms, so it cannot see calibration at all. That is why a proper score belongs on the dashboard whenever the probability itself is consumed downstream.

\redflags
\begin{itemize}
    \item Stating properness without the expectation over $y \sim q$---the definition is about expected score under the true distribution.
    \item Asserting ``accuracy is not a proper scoring rule'' without the precision that it is proper but not strictly proper; the failure that matters is the tie, not impropriety.
    \item Claiming a proper score guarantees a good model. A constant base-rate predictor is perfectly calibrated; the refinement term is what separates useful from useless.
    \item Confusing the two directions of KL.
\end{itemize}

\interviewtest{Whether you can move from a definition to a two-line proof, and whether you can say what properness buys operationally. The decomposition is the discriminator---it is the reason the concept exists rather than a technicality.}

\goodgreat
\begin{itemize}
    \item \badgeLfive{L5} Correct definition and the first-derivative argument that $p=q$ is the minimizer.
    \item \badgeLsix{L6} Gives the KL identity, does Brier in parallel, and explains why properness implies an incentive for calibrated probabilities.
    \item \badgeLseven{L7} Gives the calibration--refinement decomposition and uses it to explain why a proper score cannot be gamed, contrasts Brier's bounded penalty with log loss's unbounded one under label noise (one mislabeled row at $p=10^{-6}$ contributes 13.8 nats), and notes that Brier's uncertainty term makes raw scores incomparable across datasets with different base rates.
\end{itemize}

\followups{``Show me Brier.'' ``Which would you optimize if 2\% of labels are wrong?'' ``Is ECE a proper scoring rule?'' ``Why can accuracy not be used to select among probability models?''}
\end{interviewq}

\begin{interviewq}
\textbf{Q: Which classical models give you calibrated probabilities out of the box, and which do not?}

\badgeLfive{L5} \quad \badgetype{Conceptual} \quad \badgefreq{\freqmed}

\stronganswer

Organize by \textit{mechanism}, not by a list. A model tends to be calibrated when it is trained on a strictly proper scoring rule with the matching link and is not heavily regularized; it tends to be miscalibrated when the training objective optimizes something else (margins), when a structural assumption is violated in a way that compounds evidence, or when the training base rate is not the serving base rate.

\textbf{Logistic regression: good, by construction with fine print.} The MLE stationarity condition is $X^\top(y - \hat p) = 0$; with an intercept column this forces the mean predicted probability to equal the observed base rate and zero residual correlation with every feature, so it is calibrated in-sample along every direction the features span. That is an in-sample property: heavy regularization, a misspecified link, or distribution shift all break it.

\textbf{Naive Bayes: badly overconfident.} The conditional-independence assumption double-counts correlated evidence---ten near-duplicate features vote ten times---so posteriors are driven toward 0 and 1. The ranking often survives (which is why NB works as a classifier) while the magnitudes are meaningless. This is the canonical ``accurate but badly calibrated'' model.

\textbf{SVM: no probabilities at all.} Hinge loss produces a margin, and distance to a hyperplane is not a likelihood. Platt scaling was invented for exactly this.

\textbf{Random forest: under-confident, compressed.} The output is a vote fraction, an average of $\{0,1\}$ votes, and a forest is essentially never unanimous, so predictions bunch away from the extremes---the reliability curve is characteristically S-shaped in the opposite direction from a boosted model's.

\textbf{Boosted trees: fair, drifting overconfident.} Log-loss boosting starts from a proper objective, which helps a great deal, but continued rounds push margins outward; check calibration after early stopping. AdaBoost, whose exponential loss is proper for a different link, needs an explicit sigmoid map before its scores mean anything.

\textbf{kNN: coarse.} Probabilities are $j/k$, so $k = 10$ has eleven expressible values and cannot say 0.05.

\textbf{Any model trained on rebalanced data: systematically shifted}, regardless of family---the fix is the odds correction, not a calibrator.

Close with the operational rule: calibration is cheap to check (reliability diagram plus a proper score on held-out data) and cheap to fix (a monotone map on a held-out calibration set), so the answer is never ``this family is fine, skip it''---it is ``measure it, on the slices you make decisions about.''

\redflags
\begin{itemize}
    \item ``Logistic regression is always calibrated''---true in-sample, false under regularization, misspecification, or shift.
    \item Believing random forests are overconfident (they are the opposite) or that SVM decision values are probabilities.
    \item Offering no mechanism, only a memorized list.
    \item Forgetting that resampling for class imbalance destroys calibration in every family.
\end{itemize}

\interviewtest{Whether you can connect a training objective to the shape of its reliability curve. This is the classical-ML value-add on a topic whose mechanics are otherwise generic---and the follow-up is always ``so what would you do about it.''}

\goodgreat
\begin{itemize}
    \item \badgeLfive{L5} Correct classification of four or five families with a reason for each.
    \item \badgeLsix{L6} Gives the score-equation argument for logistic regression with its fine print, the independence-violation mechanism for naive Bayes, the vote-fraction compression for RF, and the downsampling correction formula.
    \item \badgeLseven{L7} Notes that ranking often survives miscalibration so the model may still be fine for retrieval, that marginal calibration does not imply subgroup calibration, and that the calibration map itself is a model with a bias--variance trade (Platt vs.\ isotonic by calibration-set size).
\end{itemize}

\followups{``Why is a random forest under-confident specifically?'' ``You downsampled negatives 10:1---what happens and how do you undo it?'' ``How much held-out data does isotonic need?''}
\end{interviewq}

\begin{interviewq}
\textbf{Q: Confidence interval versus prediction interval---define both, and tell me which one the product manager actually wanted.}

\badgeLfive{L5} \quad \badgetype{Conceptual} \quad \badgefreq{\freqmed}

\stronganswer

A \textbf{confidence interval} covers a parameter or a conditional \textit{mean}: ``the average delivery time on this route is 31--33 minutes.'' A \textbf{prediction interval} covers a \textit{new observation}: ``the next delivery on this route takes 18--46 minutes.'' In OLS at a test point,
\begin{equation*}
\text{CI: } \hat y_0 \pm t\,\hat\sigma\sqrt{x_0^\top(X^\top X)^{-1}x_0}, \qquad
\text{PI: } \hat y_0 \pm t\,\hat\sigma\sqrt{1 + x_0^\top(X^\top X)^{-1}x_0},
\end{equation*}
and the entire difference is the $1$ under the second radical---the new observation's own noise. The consequence is the memorable part: as $n$ grows the CI half-width goes to zero while the PI half-width converges to $t\hat\sigma$. With $\sigma = 10$ minutes and $n = 400$, the 95\% CI half-width is $1.96 \times 10/\sqrt{400} = 0.98$ minutes and the PI half-width is $1.96 \times 10\sqrt{1+1/400} = 19.6$ minutes---twenty times wider. \textbf{A prediction interval never shrinks below the noise floor.}

Map it back to the chapter's vocabulary: the CI captures epistemic uncertainty, which shrinks with data; the PI adds aleatoric uncertainty, which does not. In the conjugate Gaussian model this is literally the posterior predictive variance $\tau_n^2 + \sigma^2$, with $\tau_n^2 \to 0$ and $\sigma^2$ fixed.

And the product manager almost always wanted the prediction interval. ``How late might \textit{my} package be'' is a question about one draw, not about a mean, and quoting the CI understates the answer by that factor of twenty---which is exactly how teams end up shipping confident promises they miss constantly. Add the honest caveat that both formulas assume the model is right; if the mean function is misspecified, neither interval has any protection, which is one motivation for the distribution-free machinery of the conformal question.

\redflags
\begin{itemize}
    \item Using the terms interchangeably, or saying the PI is ``the CI with a bigger $t$.''
    \item Claiming a prediction interval shrinks to zero with enough data.
    \item Not noticing the two intervals answer different questions and only one of them is the one that was asked.
\end{itemize}

\interviewtest{A precision check on a distinction that separates people who have used statistics from people who have read about it, and a direct probe of the aleatoric/epistemic split.}

\goodgreat
\begin{itemize}
    \item \badgeLfive{L5} Correct definitions and the ``PI is wider'' statement with the reason.
    \item \badgeLsix{L6} Writes both formulas, works the numbers, and maps them onto aleatoric/epistemic and the posterior predictive variance.
    \item \badgeLseven{L7} Notes that both are model-dependent and that misspecification voids them, that heteroscedasticity means a constant-width PI is wrong in the regions that matter, and reaches for quantile regression or conformal prediction when the assumptions are unattractive.
\end{itemize}

\followups{``How do you get a prediction interval from a gradient-boosted tree?'' ``What if the noise varies with $x$?'' ``How would you validate that your 90\% interval really covers 90\%?''}
\end{interviewq}

\begin{interviewq}
\textbf{Q: I want prediction sets with a coverage guarantee and no distributional assumptions. Build it, then tell me when the guarantee is worthless.}

\badgeLsix{L6} \quad \badgetype{First Principles} \quad \badgefreq{\freqmed}

\stronganswer

Split conformal prediction. \textbf{Recipe:} hold out a calibration set of size $n$, disjoint from training; fit $\hat f$ on the training set only and freeze it; compute nonconformity scores on the calibration set---$s_i = |y_i - \hat f(x_i)|$ for regression, $s_i = 1 - \hat p_{y_i}(x_i)$ for classification; take $\hat q$ to be the $\lceil (n+1)(1-\alpha)\rceil/n$ empirical quantile of those scores; and output $C(x) = \{y : s(x,y) \le \hat q\}$, which is $\hat f(x) \pm \hat q$ for regression and a variable-size label set for classification. With $n = 1000$ and $\alpha = 0.1$ that means the 901st smallest residual, the 90.1th percentile---the $(n+1)$ is the finite-sample correction, not a rounding detail, and feasibility needs $n \ge 1/\alpha - 1$.

\textbf{Why it works, in one line:} if the calibration points and the test point are exchangeable, the rank of $s_{n+1}$ among all $n+1$ scores is uniform, so $\prob{s_{n+1} \le s_{(k)}} \ge k/(n+1)$; choosing $k = \lceil (n+1)(1-\alpha)\rceil$ gives $\prob{Y_{n+1} \in C(X_{n+1})} \ge 1-\alpha$, with an upper bound of $1-\alpha+\tfrac{1}{n+1}$ under continuous scores. Finite-sample, distribution-free, model-agnostic---the model can be anything, including a black box you did not train.

\textbf{Now the fine print, which is what the question is really about.} (1) The guarantee is \textbf{marginal, not conditional}: 90\% averaged over inputs is routinely 98\% on easy cases and 65\% on a hard subgroup, and distribution-free conditional coverage is provably impossible with non-trivial sets. If you need per-group validity, calibrate within groups (Mondrian conformal); if you want widths that adapt to difficulty, conformalize a quantile-regression interval (CQR), which keeps the guarantee and inherits the adaptivity. (2) It says \textbf{nothing about width}: a useless model gets exact coverage with enormous intervals---coverage is validity, efficiency comes only from the base model. (3) The coverage is also marginal \textbf{over the calibration draw}: for a fixed calibration set the realized coverage is Beta$(n+1-\ell,\ell)$ with $\ell = \lfloor(n+1)\alpha\rfloor$, so at $n = 100$, $\alpha = 0.1$ it has mean 0.901 and standard deviation 3.0 points---your deployment could really be at 85\%---dropping to 0.95 points at $n = 1000$. (4) \textbf{Exchangeability is the whole ballgame.} Under temporal drift, feedback loops, or a train/serve mismatch, the guarantee is void with no error message. Weighted conformal restores it under covariate shift if you can estimate the likelihood ratio; adaptive conformal inference adjusts $\alpha_t$ online for streaming data, trading finite-sample validity for long-run average coverage; and in every case you recalibrate on a rolling window and \textbf{monitor realized coverage as a production metric}.

\redflags
\begin{itemize}
    \item Claiming conditional (per-input or per-group) coverage. This is the most common overstatement and the follow-up is guaranteed.
    \item Fitting the model on the calibration data---the guarantee dies immediately.
    \item Using the plain $(1-\alpha)$ quantile instead of $\lceil(n+1)(1-\alpha)\rceil/n$, or not handling the small-$n$ infeasible case.
    \item Presenting it as robust to distribution shift. It is distribution-\textit{free}, which is a different property: it makes no assumption about \textit{which} distribution, but requires that it not change.
    \item Treating it as a way to improve the model.
\end{itemize}

\interviewtest{Whether a 2020s technique you name-drop is one you actually understand. The mechanics are five lines; the discriminator is the precise account of what the guarantee covers---marginal over inputs and over the calibration draw, and conditional on exchangeability.}

\goodgreat
\begin{itemize}
    \item \badgeLfive{L5} Correct split-conformal recipe and the statement of marginal, distribution-free, finite-sample coverage.
    \item \badgeLsix{L6} Gives the exchangeability/rank-uniformity argument, the $(n+1)$ correction with a worked quantile index, and the marginal-versus-conditional distinction with Mondrian or CQR as the remedy.
    \item \badgeLseven{L7} Adds the Beta distribution of realized coverage and sizes the calibration set from it, separates validity from efficiency, and lays out the shift story concretely---weighted conformal under known covariate shift, adaptive conformal for streaming, and coverage monitoring as the operational backstop.
\end{itemize}

\followups{``Your sets are technically valid but huge---what do you do?'' ``How would you get 90\% coverage \textit{within each country}?'' ``The data is a time series---what breaks?'' ``How large should the calibration set be?''}
\end{interviewq}

\begin{interviewq}
\textbf{Q: Logistics wants a P90 delivery-time estimate per order, not a point prediction. What do you build, and why is your loss function correct?}

\badgeLsix{L6} \quad \badgetype{Trade-off} \quad \badgefreq{\freqlow}

\stronganswer

Start by pinning the ask: P90 means a conditional quantile, $q_{0.9}(x)$ such that 90\% of actual times for orders like this fall below it. That is not the mean plus a fudge factor, and it is not the mean plus $1.28\sigma$ unless the conditional distribution is Gaussian---delivery times are right-skewed, so that shortcut under-promises in the tail exactly where the promise matters.

\textbf{Build:} quantile regression, or a GBDT with the pinball objective at $\tau = 0.9$. The loss is
\begin{equation*}
\rho_\tau(u) = u\big(\tau - \mathds{1}[u<0]\big), \qquad u = y - q,
\end{equation*}
i.e.\ asymmetric absolute error weighting under-prediction $\tau$ and over-prediction $1-\tau$. \textbf{Why it is correct}, derived: with $R(q) = \tau\int_q^\infty (y-q)dF + (1-\tau)\int_{-\infty}^q (q-y)dF$, differentiating gives $R'(q) = -\tau(1-F(q)) + (1-\tau)F(q) = F(q) - \tau$, so $R'(q) = 0$ exactly at $F(q) = \tau$, and $R''(q) = f(q) \ge 0$ makes it the global minimum. The minimizer of the pinball loss \textit{is} the $\tau$-quantile. Sanity check: $\tau = 0.5$ reduces to $\tfrac{1}{2}|u|$ and gives the median, the same way squared error gives the mean.

\textbf{The trade-off against the alternatives.} A \textit{bootstrap} interval on the prediction is the wrong tool---it captures epistemic uncertainty about the fitted function and mostly not the aleatoric spread of individual deliveries, which is what a P90 promise is about; it also costs $B$ refits. A \textit{Gaussian variance head} is cheaper but imposes symmetry on a skewed target. \textit{Quantile regression} is the right primary model: one head per $\tau$, widths that vary with $x$ (rush hour and dense urban routes get wider), and no distributional assumption. Its weakness is that it comes with \textbf{no coverage guarantee}---the fitted 90th percentile covers 90\% only if the model is right---so the strong answer wraps it in \textit{conformalized quantile regression}: calibrate the quantile interval on held-out data to obtain finite-sample marginal coverage while keeping the adaptive widths.

Practical caveats to raise unprompted: separately fitted quantiles can cross, fixed by joint fitting, post-hoc sorting, or monotone constraints; the pinball loss has zero second derivative, so second-order boosters fall back to unit Hessians or per-leaf line search (Table~\ref{cml:tab:boosting_losses}); and the business metric is realized coverage per segment, so instrument it---a P90 that is 97\% accurate in suburbs and 71\% downtown is failing the users who complain, which is the marginal-versus-conditional issue again.

\redflags
\begin{itemize}
    \item Predicting the mean and adding $1.28$ standard deviations without checking skew.
    \item Proposing MSE and then reading a quantile off the residuals as if the model had targeted it.
    \item Unable to say why the pinball loss yields a quantile---the derivation is the point of the question.
    \item Bootstrapping the model and calling the result a prediction interval, conflating epistemic and aleatoric uncertainty.
\end{itemize}

\interviewtest{Whether you can derive a loss from the estimand instead of picking one from a menu---the same instinct as the noise-model table earlier in this chapter---and whether you know that a nominal quantile is not a guaranteed one.}

\goodgreat
\begin{itemize}
    \item \badgeLfive{L5} Names quantile regression and the pinball loss, with the correct asymmetry.
    \item \badgeLsix{L6} Derives $R'(q) = F(q)-\tau$ and the median special case, explains why the mean-plus-$z\sigma$ shortcut fails on skewed targets, and raises quantile crossing.
    \item \badgeLseven{L7} Distinguishes aleatoric from epistemic to reject the bootstrap answer, wraps the model in conformalized quantile regression for finite-sample coverage, and defines realized coverage per segment as the operating metric.
\end{itemize}

\followups{``Now derive that MSE gives the mean and MAE gives the median.'' ``Two of your quantile curves cross---fix it.'' ``How would you know in production whether your P90 is really a P90?''}
\end{interviewq}

\begin{interviewq}
\textbf{Q: We are launching an automated credit-limit model. Design the uncertainty layer: what we estimate, how, and what we monitor.}

\badgeLseven{L7} \quad \badgetype{System Design} \quad \badgefreq{\freqmed}

\stronganswer

Start from the decisions, not the methods. Limits are set by an expected-value calculation over a predicted default probability, and there are three distinct consumers: the automated decision (needs a calibrated probability and a cost-derived threshold), the human review queue (needs to know \textit{which} cases the model is unsure about), and the regulator or risk committee (needs an auditable interval and a stated failure mode). Each wants a different quantity.

\textbf{1. Separate the two uncertainties.} Aleatoric---two indistinguishable applicants, one defaults---does not shrink with more rows and is what the calibrated probability itself encodes. Epistemic---thin-file applicants, new segments, a product launched last quarter---does shrink with data and is what should route a case to human review. Conflating them produces the classic failure: a confident 3\% default estimate for a segment the model has seen forty times.

\textbf{2. Estimate them.} Aleatoric: train on a strictly proper scoring rule, then recalibrate on a held-out, recent window (Platt under $\sim$1K points, isotonic above $\sim$10K), and re-derive thresholds from the cost ratio $p^\star = c_{\text{FP}}/(c_{\text{FP}}+c_{\text{FN}})$ rather than from round numbers. Epistemic: bootstrap or ensemble disagreement across resampled fits---cheap, model-agnostic, and directly usable as a routing signal; for per-segment rates with few observations, a Beta-Binomial/empirical-Bayes shrinkage layer so a merchant or ZIP with twelve observations is pulled toward the population prior instead of reporting $0/12 = 0\%$. \textbf{3. Wrap it.} Split conformal on top of the frozen model gives finite-sample marginal coverage for the reported interval, which is the thing an auditor can be told without a distributional assumption---and calibrate \textit{within} protected or business-critical groups (Mondrian) because marginal coverage will otherwise be borrowed from the easy majority.

\textbf{4. Monitor.} Four production metrics, not one: a proper score (log loss or Brier) rather than AUC alone; calibration by slice on a rolling window, since calibration degrades under drift faster than discrimination does; realized conformal coverage, which is the alarm for broken exchangeability; and the review-queue rate, since a rising epistemic signal is the earliest indicator of a population shift. Recalibrate on a schedule, not on incident.

\textbf{5. State what none of this protects against.} Every method here quantifies noise and estimation error; none quantifies \textbf{bias from misspecification or from a missing feature}, and none survives a genuine regime change---a conformal interval built on pre-recession data says nothing about a recession. That honesty, plus a documented human-override path, is what makes the design defensible rather than merely instrumented.

\redflags
\begin{itemize}
    \item One number called ``confidence'' serving all three consumers.
    \item Claiming more data will fix aleatoric uncertainty, or that conformal prediction handles distribution shift.
    \item No monitoring: uncertainty estimates computed once at training and never validated against realized outcomes.
    \item Reporting $0/12 = 0\%$ default for a thin segment because the MLE said so.
    \item Presenting the uncertainty layer as protection against model bias, which it is not.
\end{itemize}

\interviewtest{Whether you can assemble the chapter's pieces into a system with named consumers, chosen methods, and monitoring---and whether you will state the limits of your own design without being pushed.}

\goodgreat
\begin{itemize}
    \item \badgeLfive{L5} Names calibration plus one uncertainty method and proposes monitoring.
    \item \badgeLsix{L6} Separates aleatoric from epistemic with a method for each, derives the threshold from costs, and monitors calibration by slice on a rolling window.
    \item \badgeLseven{L7} Maps each quantity to the decision that consumes it, adds empirical-Bayes shrinkage for thin segments and group-conditional conformal for coverage that cannot be borrowed from the majority, treats realized coverage as the drift alarm, and states plainly that none of it bounds misspecification bias.
\end{itemize}

\followups{``An applicant's interval is wide---what does the system actually do?'' ``How do you validate the uncertainty layer before launch?'' ``Half your calibration data is now stale---what do you drop?'' ``What would you tell a regulator the guarantee means?''}
\end{interviewq}
